% Historical appendix fragment before consolidation; not a standalone document.
% Its table labels and discussion are preserved for traceability, not inclusion.
\appendix

% DATA FIGURE PROVENANCE: {"dataset": "Dataset203_GeoFWI3D", "image": "datasets/Dataset203_GeoFWI3D/imagesTr/model_0001_0000.nii.gz", "labels": "datasets/Dataset203_GeoFWI3D/labelsTr/model_0001.nii.gz", "shape": [96, 96, 96], "plane_axis": 0, "slice_index_zero_based": 50, "display_gray_limits": [-1.0, 1.0], "display_gray_levels": 32, "display_flip_axis2": false, "labels_in_slice": [0, 1], "image_sha256": "f2a8f1a7e2ac61692301b6780c5a7255f98cfa50ca7017aa7b7f0582c858921a", "labels_sha256": "13cbf533c902d76848279d8636a330f64bb48d51371cecf546046d03f530eb33"}
% DATA FIGURE PROVENANCE: {"dataset": "Dataset101_MSD", "image": "datasets/Dataset101_MSD/imagesTr/hippocampus_001_0000.nii.gz", "labels": "datasets/Dataset101_MSD/labelsTr/hippocampus_001.nii.gz", "shape": [35, 51, 35], "plane_axis": 0, "slice_index_zero_based": 14, "display_gray_limits": [9.0, 106.0], "display_gray_levels": 32, "display_flip_axis2": true, "labels_in_slice": [0, 1, 2], "image_sha256": "ab47d6115b412ffaba041cbf48962967127f199f39eb56c777a6bb3746353b4a", "labels_sha256": "31e218cd74ee787c25a17a9cc3a13743ba8ecd11dd76fbad3c0c18c09fd8b396"}
% DATA FIGURE PROVENANCE: {"dataset": "Dataset103_ADNI", "image": "datasets/Dataset103_ADNI/imagesTr/hippocampus_adni_10007_L_0000.nii.gz", "labels": "datasets/Dataset103_ADNI/labelsTr/hippocampus_adni_10007_L.nii.gz", "shape": [32, 48, 40], "plane_axis": 0, "slice_index_zero_based": 15, "display_gray_limits": [75.62541742107685, 787.0265820841897], "display_gray_levels": 32, "display_flip_axis2": true, "labels_in_slice": [0, 1], "image_sha256": "b1479c7bd284e5d4013dd970ff5330819c49732d8f3b8b27c5d4957bde68b799", "labels_sha256": "03fbe773f988e7541ced3681d1f76719745059b8d404b9c9de76c74401def8b3"}
% CURRENT NUMERICAL PROVENANCE -- 21 October 2026:
% reports/thesis_noaug_update_20261021/audit.json records source SHA-256 identities.
% Conventional MSD and ADNI final reports: 504 runs per dataset.
% MSD MedSAM3: one campaign; execution batches preserve 27 and 108 records.
% ADNI MedSAM3: 126 complete run archives, rebuilt beyond the stale summary.
% Performance tables retain the recorded matched baselines for each comparison.
% SAM3 final 126-run aggregates verified in reports/thesis_sam3_final_20261023/.
% All conventional and adaptation objective inventories are complete.

\section{Numerical and Implementation Details}
\label{sec:technical-details}

The main constraint chapter defines the supervised properties and their
losses. This appendix supplies the stable evaluations, derivatives and
protocol-specific reductions needed to reproduce them.

\subsection{Stable Membership Evaluation}

For a class group $F$ and its complement $\overline F$, suppress the
case and voxel indices. For any class set $U$, $z_U$ collects its logits
and the log-sum-exp operation is
$\operatorname{LSE}(z_U)=\log\sum_{c\in U}\exp z_c$.
The sigmoid $\sigma(s)=1/(1+\exp(-s))$ converts a binary logit into a
membership probability. The grouped logit
\begin{equation}
s_F=\operatorname{LSE}(z_F)-\operatorname{LSE}(z_{\overline F})
\label{eq:grouped-predicate-logit}
\end{equation}
satisfies $\sigma(s_F)=\sum_{c\in F}p_c$. On MSD, $F=\{A,P\}$
and the complement contains only background. Restoring case and voxel
indices, write this grouped logit as $s_i(v)$. With
$\operatorname{softplus}(s)=\log(1+\exp s)$, the stable membership loss is
\begin{equation}
e_i(v)=h_i(v)\operatorname{softplus}(-s_i(v))
       +(1-h_i(v))\operatorname{softplus}(s_i(v)).
\label{eq:band-stable-bce}
\end{equation}
Grouped logits and the combined logit-based loss use float32.
Classwise log-probabilities use log-softmax,
$\log p_c=z_c-\operatorname{LSE}(z_{\mathcal C})$. These evaluations avoid taking logarithms of rounded
probabilities.

Adding $\varepsilon$ to a nonempty support count multiplies its exact
mean by $\rho=|U|/(|U|+\varepsilon)$. Thus, for band-side means
$b_i^s=-\ell_\theta(\Phi_i^s)$ and
$\rho_i^s=n_i^s/(n_i^s+\varepsilon)$,
$\mathcal L_{B,i}=\tfrac12\sum_s\rho_i^s b_i^s$ and
\begin{equation}
\exp(-\mathcal{L}_{B,i})
=\mathcal{G}_\theta(\Phi_i^{\mathrm{in}})^{\rho_i^{\mathrm{in}}/2}
 \mathcal{G}_\theta(\Phi_i^{\mathrm{out}})^{\rho_i^{\mathrm{out}}/2}.
\label{eq:balanced-band-satisfaction}
\end{equation}
The count factors are retained in both values and gradients. Empty
supports follow the validity rules of the corresponding protocol.

\paragraph{Boundary-loss settings.}
\label{sec:band-implementation-details}

The principal MSD implementation evaluates grouped logits and binary
cross-entropy in float32. Optional focal factors and weights based on the
number of neighbours are disabled. Both full-volume bands must be nonempty
for a case to enter the mean; if no valid case remains, the code returns a
differentiable zero connected to the logits. The count constant is
$\varepsilon=10^{-6}$. Equation~\ref{eq:balanced-band-satisfaction} retains
its factor in the loss and gradient.


\subsection{Relational and Signed-Distance Derivatives}

Band and depth losses have the usual weighted cross-entropy derivatives.
For edges, let $\mathcal N_i^B(u)$ contain the neighbours of $u$ in
the selected pair graph. The case-loss derivatives are
\begin{equation}
\begin{aligned}
\frac{\partial\mathcal{L}_{E,i}}{\partial q_i(u)}
 &=\frac{2}{n_i^E}\sum_{v\in\mathcal{N}_i^B(u)}
       \bigl(\delta_i(u)-\delta_i(v)\bigr),\\
\frac{\partial\mathcal{L}_{E,i}}{\partial s_i(u)}
 &=q_i(u)\bigl(1-q_i(u)\bigr)
   \frac{\partial\mathcal{L}_{E,i}}{\partial q_i(u)}.
\end{aligned}
\label{eq:edge-logit-gradient}
\end{equation}
The first expression sums the differences between a voxel's membership
residual and those of its neighbours, an operation called the graph
Laplacian. The sigmoid factor $q_i(u)(1-q_i(u))$ becomes small near
probabilities zero and one, reducing the logit gradient.

For the MSD signed-distance term, define $\bar\varphi_{i,0}=0$ and
$\bar\varphi_{i,c}=\varphi_{i,c}$ for $c\in\{A,P\}$. Then
\begin{equation}
\frac{\partial\mathcal{L}_{\mathrm{SDF},i}}
     {\partial z_{\theta,i,k}(v)}
=\frac{p_{i,k}(v)}{2|\Omega_i|}
 \left[\bar\varphi_{i,k}(v)
       -\sum_{c\in\mathcal{C}}p_{i,c}(v)\bar\varphi_{i,c}(v)\right].
\label{eq:sdf-logit-gradient}
\end{equation}
The annotation-dependent offset that makes the loss nonnegative does
not change its derivative. Batch reduction, calibrated coefficients and
activation schedules multiply these case derivatives as specified in
Equation~\ref{eq:scheduled-constraint-objective}.

\subsection{Physical Geometry of the Depth Construction}

For the MSD hippocampal union $M_i$, the centroid $\mu_i$ is the mean
physical position and the covariance matrix $\Sigma_i$ describes how
foreground positions spread around it:
\begin{equation}
\begin{aligned}
\mu_i&=\frac{1}{|M_i|}\sum_{v\in M_i}r_i(v),\\
\Sigma_i&=\frac{1}{|M_i|}\sum_{v\in M_i}
       (r_i(v)-\mu_i)(r_i(v)-\mu_i)^{\mathsf T}.
\end{aligned}
\label{eq:depth-centroid-covariance}
\end{equation}
The unit eigenvector associated with the largest eigenvalue of
$\Sigma_i$ gives the principal axis $\ell_i$ used in the main text.
It is oriented anteriorly to choose its sign. Reversing this sign leaves the projected rightward normal in
Equation~\ref{eq:depth-plane-normal} unchanged. The gap between the largest two eigenvalues indicates how clearly one
direction dominates the spread and is recorded as a diagnostic, without replacing
ambiguous geometry by a zero constraint.

The surrounding region is
\begin{equation}
R_i=\left\{v\in\Omega_i:
       \min_{u\in M_i}\|r_i(v)-r_i(u)\|_2\leq m\right\},
\qquad m=8\,\mathrm{mm}.
\label{eq:depth-region-of-interest}
\end{equation}
with the same coordinate map as the plane. The distance transform uses
column norms of $Q_i$ as axis spacings. This implements physical
Euclidean distance for orthogonal voxel axes, including rotated grids.
Sheared affines are rejected, as are annotations with fewer than four
foreground points or a projected normal of norm below 0.1. Such
failures require correction of the preparation.

The conventional medical recipes use one $64^3$ volume per update and apply
constraints only to the primary output; nnU-Net's other supervised
resolutions retain their baseline losses. Identity affines introduced during
ADNI export leave grid coordinates unchanged; these placeholder transforms
do not establish anatomical orientation or physical spacing; its distance
and depth constructions therefore use stored array coordinates in voxel
units.

\subsection{GeoFWI3D Observation Boundaries}
\label{sec:geo-boundaries}

The final four-model geological campaign retains model-specific band
supports. For axis lengths $N_{i,k}$, define the interior domain
\begin{equation}
R_i^{(r)}=\{v\in\Omega_i:
 r\leq v_k\leq N_{i,k}-1-r\ \text{for }k=1,2,3\}.
\label{eq:geological-band-frame}
\end{equation}
UNETR uses $\widetilde B_i^s=B_i^s\cap R_i^{(2)}$, excluding the
two-voxel observation frame. nnU-Net, SwinUNETR and the compact U-Net
use $\widetilde B_i^s=B_i^s$ without this exclusion. In every recipe,
the band loss is the half-sum of inner and outer binary cross-entropy (BCE) sums divided
by their respective counts plus $10^{-6}$. Both sides must be
nonempty; invalid cases are skipped and an entirely invalid batch
returns a differentiable zero.

Edge pairs require both endpoints in
$\widetilde B_i^{\mathrm{in}}\cup\widetilde B_i^{\mathrm{out}}$.
The reduction is an exact pair mean followed by a valid-case mean,
with differentiable zero when no valid case remains. Thus only the
UNETR recipe excludes pairs involving the observation frame.
The single-foreground distance loss is
\begin{equation}
\mathcal{L}_{\mathrm{SDF},i}^{\mathrm{Geo}}
=\frac{1}{|\Omega_i|}\sum_{v\in\Omega_i}
 \bigl(p_{i,S}(v)-y_{i,S}(v)\bigr)\varphi_{i,S}(v).
\label{eq:geological-sdf-loss}
\end{equation}
It covers the full stored grid, including the observation frame,
and retains the half-voxel shift, unit-grid Euclidean distances and
clipping threshold eight. The coordinate and depth conventions are
specified in Section~\ref{sec:geo-methods}.

\subsection{Sampled Volumetric Supervision in Pretrained Models}
\label{sec:sampled-pretrained-supervision}

Native three-dimensional class masks define separate bands for
$c\in\{0,A,P\}$ on MSD, $c=H$ on ADNI and $c=S$ for SAM3 on
GeoFWI3D. Zero padding is retained, including for MSD background,
whose inner band can therefore reach crop faces. Support counts are
computed over the full volume before sampling.

Let $Z_i$ be the number of planes. One plane, the anchor, is selected
uniformly with index $k$, and $\Omega_i^{(k)}$ denotes its voxel domain.
For a class with two nonempty full-volume band sides, its sampled
contribution is
\begin{equation}
\begin{aligned}
\widehat{\mathcal{L}}_{B,i,c}^{(k)}
=-\frac{Z_i}{2}\Biggl[
 &\frac{\sum_{v\in B_{i,c}^{\mathrm{in}}\cap\Omega_i^{(k)}}
                  \log p_{\theta,i,c}(v)}
        {n_{i,c}^{\mathrm{in}}+\varepsilon}\\
 {}+{}&\frac{\sum_{v\in B_{i,c}^{\mathrm{out}}\cap\Omega_i^{(k)}}
                  \log(1-p_{\theta,i,c}(v))}
        {n_{i,c}^{\mathrm{out}}+\varepsilon}\Biggr].
\end{aligned}
\label{eq:medsam-sampled-band-loss}
\end{equation}
where $n_{i,c}^s=|B_{i,c}^s|$ and $\varepsilon=10^{-6}$.
For fixed predictions, averaging over anchors recovers the full-volume
stabilized loss. An empty full-volume side gives a zero class contribution;
an empty sampled support also contributes zero without changing the
global denominator.

Let $\mathcal E_{i,c}^B(k)$ contain selected in-plane pairs on $k$
and through-plane pairs from $k$ to $k+1$ for $k<Z_i-1$.
This assigns every full-volume edge to exactly one anchor. With
$n_{i,c}^E=|\mathcal E_{i,c}^B|>0$ and
$\delta_{i,c}=p_{\theta,i,c}-y_{i,c}$,
\begin{equation}
\widehat{\mathcal{L}}_{E,i,c}^{(k)}
=\frac{Z_i}{n_{i,c}^E}
  \sum_{\{u,v\}\in\mathcal{E}_{i,c}^B(k)}
    \bigl(\delta_{i,c}(u)-\delta_{i,c}(v)\bigr)^2.
\label{eq:medsam-sampled-edge-loss}
\end{equation}
satisfies
\begin{equation}
\frac{1}{Z_i}\sum_{k=0}^{Z_i-1}
  \widehat{\mathcal{L}}_{E,i,c}^{(k)}
=\frac{1}{n_{i,c}^E}\sum_{\{u,v\}\in\mathcal{E}_{i,c}^B}
    \bigl(\delta_{i,c}(u)-\delta_{i,c}(v)\bigr)^2
=\mathcal{L}_{E,i,c}.
\label{eq:medsam-edge-expectation}
\end{equation}
The native supervised objective uses companion plane
$(k+1)\bmod Z_i$, where the modulo operation wraps the last index back
to zero. The edge estimator never joins the final plane to the first. Empty full-volume edge support contributes zero.

On MSD, both estimators average all three classes with weight $1/3$,
including zero-contributing classes. ADNI MedSAM3 and GeoFWI3D SAM3
evaluate their single foreground mask, with no additional
background-class average. Volumes and anchors are sampled
uniformly, and the two pairs per update are averaged through the
fixed accumulation schedule. For fixed slice predictions, the finite
expectation identities also apply to their gradients. They provide
volumetric supervision while retaining an image-based encoder.


Signed-distance and depth supervision also admit sampled volumetric
reductions. Let $C_+$ be the foreground class set, let each plane
contain $N_i$ voxels and let $w_i(v)$ be full-volume normalized depth
weights. Their anchor estimators are
\begin{align}
\widehat{\mathcal L}_{S,i}^{(k)}
 &=\frac{1}{|C_+|N_i}\sum_{c\in C_+}
   \sum_{v\in\Omega_i^{(k)}}(p_{i,c}(v)-y_{i,c}(v))\varphi_{i,c}(v),
   \label{eq:medsam-sampled-sdf}\\
\widehat{\mathcal L}_{D,i}^{(k)}
 &=-Z_i\sum_{v\in\Omega_i^{(k)}}w_i(v)\log p_{i,y_i(v)}(v).
   \label{eq:medsam-sampled-depth}
\end{align}
Uniform anchor averaging recovers the full-volume SDF mean and weighted
cross-entropy. The $Z_i$ correction is needed for the full-volume depth
weights, whereas SDF already averages within each plane. Distances use the
verified 1 mm MSD grid and unit source-array coordinates for ADNI and
GeoFWI3D SAM3, as specified in the methods. The latter two constructions
do not infer physical millimetres or anatomical orientation from their
working coordinate maps.

\subsection{Calibration and Discrete Compute Limits}

Sequential calibration adds weighted gradient vectors before measuring
their root mean square (RMS). Let $g_{i,0}$ and $g_{i,B}$ be the
baseline and band-loss gradients with respect to the $m_i$ primary
logits of case $i$. For a vector $g$ with $m_i$ entries,
$\operatorname{RMS}(g)=\sqrt{m_i^{-1}\sum_{k=1}^{m_i}g_k^2}$.
For example,
\begin{equation}
\begin{aligned}
\operatorname{RMS}(g_{i,0}+\lambda_B g_{i,B})^2
={}&\operatorname{RMS}(g_{i,0})^2
 +2\lambda_B\langle g_{i,0},g_{i,B}\rangle_{m_i}\\
 &+\lambda_B^2\operatorname{RMS}(g_{i,B})^2,
\end{aligned}
\label{eq:calibration-cross-term}
\end{equation}
where $\langle a,b\rangle_{m_i}=m_i^{-1}\sum_k a_kb_k$.
The cross term retains directional agreement. Adding norms instead
would generally give another coefficient. Quantiles use linear
interpolation, and calibration requires finite positive ratios,
without an arbitrary fallback for an invalid result.

Let $d_a$ and $d_a^S$ be the update boundaries for early and delayed
SDF activation. The measured per-update costs $f_{a,0}$, $f_{a,1}$
and $f_{a,2}$ correspond to baseline-only training, early terms active,
and delayed SDF also active. The cumulative charge for $n$ updates is
\begin{equation}
\begin{aligned}
C_a(n)={}&f_{a,0}\min(n,d_a)\\
 &+f_{a,1}\bigl[\min(n,d_a^S)-d_a\bigr]_+\\
 &+f_{a,2}[n-d_a^S]_+,
\end{aligned}
\label{eq:piecewise-optimization-cost}
\end{equation}
Here $[x]_+=\max(x,0)$ and $C_{\max}$ denotes the medical optimization
ceiling specified in Equation~\ref{eq:main-numerical-cap}. The largest
integer update count satisfying that ceiling is permitted. The delayed boundary is chosen nearest
to charge $5C_{\max}/9$, after the early ramp, with ties resolved
toward the earlier boundary. Shared activation requires only two
cost phases.

Attempted updates are charged even when a numerical failure prevents
a parameter change; successful updates are tracked separately.
Fractional final updates are not permitted. Calibration remains
outside this optimization ledger.

\clearpage

\subsection{Detailed Calibration Protocol}\label{sec:detailed-calibration}

Loss magnitudes depend on their supports and reductions. Calibration
therefore matches measured gradient scales before assigning fixed
coefficients. It uses the ten support cases after nine
baseline updates for nnU-Net and 50 for the other models. The network
is evaluated without changing its parameters, other stored network
values or random-generator states. The float32 primary logits $z_i$, with $m_i$ entries, are detached
from the network: derivatives are measured with respect to these outputs
without propagating back to its parameters. Define
\begin{equation}
g_{i,0}=\nabla_{z_i}\mathcal{L}_{\mathrm{seg},i},
\qquad
g_{i,j}=\nabla_{z_i}\mathcal{L}_{j,i},
\qquad
\operatorname{RMS}(g)=
\sqrt{\frac{1}{m_i}\sum_{k=1}^{m_i}g_k^2}.
\label{eq:calibration-logit-gradients}
\end{equation}
nnU-Net's baseline derivative includes the primary output's native
deep-supervision weight; detached coarser outputs make no contribution
to this derivative.

Families are calibrated in the order $B\to E\to S\to D$, omitting
absent terms. If $K_j$ contains the families already calibrated, the
reference includes their weighted gradients:
\begin{equation}
g_{i,<j}=g_{i,0}+\sum_{k\in K_j}\lambda_k g_{i,k},
\qquad
r_{i,j}=\frac{\operatorname{RMS}(g_{i,j})}
                  {\operatorname{RMS}(g_{i,<j})}.
\label{eq:sequential-calibration-ratio}
\end{equation}
Let $\eta_a$ be the target median gradient ratio and $\kappa_a$ its
upper-percentile limit; $Q_{0.95}$ denotes the 95th percentile across
calibration cases. The coefficient is then
\begin{equation}
\lambda_j=
\min\left\{
\frac{\eta_a}{\operatorname{median}_i r_{i,j}},\,
\frac{\kappa_a}{Q_{0.95}(\{r_{i,j}\}_i)}
\right\}.
\label{eq:constraint-calibrated-coefficient}
\end{equation}
For nnU-Net, $(\eta_a,\kappa_a)=(0.6,0.75)$ in both datasets;
for the other three models, $(\eta_a,\kappa_a)=(0.1,0.5)$. The explicit
numeric protocol fields and saved calibration determine these settings.
The first bound targets the median gradient ratio and the second
limits its 95th percentile on the calibration cases. Adding gradient vectors
before taking their RMS retains interactions between their directions,
making the coefficient specific to the combination. Depth is calibrated
against preceding terms at full coefficients even when SDF is activated
later. These measurements use no validation labels and do not guarantee
the same gradient ratios throughout optimization.

For a baseline prefix of $d_a$ updates and a ramp lasting $w_a$ updates,
early terms follow
\begin{equation}
s_a^{\mathrm{early}}(t)=
\begin{cases}
0, & t\leq d_a,\\
\min\{1,(t-d_a)/w_a\}, & t>d_a.
\end{cases}
\label{eq:constraint-early-ramp}
\end{equation}
with baseline prefix and ramp length $d_a=w_a=9$ for nnU-Net and
$d_a=w_a=50$ otherwise. SDF either shares this schedule or starts an
analogous ramp after a delayed boundary near five ninths of the
optimization allowance.


The activation policy was developed using earlier MSD validation
experiments and then retained in these campaigns. It is part of the
configured training procedure rather than an independently selected rule. SDF
combinations containing depth were evaluated only with delayed activation.
Appendix~\ref{sec:technical-details} gives the discrete boundary and
cost calculation.

\subsection{Medical Resource Accounting}\label{sec:detailed-resource-accounting}

The medical budget was anchored to nnU-Net because its measured training
update was the most expensive among the four implementations. Exploratory
learning curves informed the choice of 90 nnU-Net updates as the minimum
practical reference allowance for this short-budget comparison. With ten
training crops and one crop per update, this represents nine complete
passes through the selected training subset. The intention was to leave
all architectures a useful learning window while keeping computation
limited. This choice concerns computational cost, rather than model size
alone: UNETR has more parameters than nnU-Net, but cheaper updates in
these configurations.

The measured native nnU-Net baseline update costs
$355{,}293{,}896{,}301$ FLOPs. Multiplying it by the 90-update reference
schedule gives the common optimization ceiling used on both MSD and ADNI:
% CAP PROVENANCE: reports/four_models_six_optimization_31976_20261006/
% prepared_plan.json, budget_origin: 90 times measured native nnU-Net
% baseline update cost, selected following exploratory learning curves.
\begin{equation}
 C_{\max}=31{,}976{,}450{,}667{,}090\ \mathrm{FLOPs}
          =31{,}976.45066709\ \mathrm{GFLOPs}.
\label{eq:main-numerical-cap}
\end{equation}
Each other architecture receives this same FLOP allowance, with its
number of updates determined by its own measured cost. Cheaper models
therefore receive more updates and complete more epochs, where an epoch
is one pass through the selected training subset. The cap matches
computational allowance across architectures; it does not require them
to reach the same stage of learning or convergence.

The arithmetic ledger charges forward computation, active losses,
backward computation and the optimizer update, counting two FLOPs per
multiply--accumulate. An update is admitted only if its charge fits
within the remaining allowance. This defines the optimization resource
being compared; it does not measure GPU utilization or elapsed time.

Constraint calculations in the optimization graph are included.
Construction of fixed reference targets, validation,
loading and storage are excluded. Calibration is accounted separately.
Thus, the matched quantity is optimization computation, while total
training-plus-calibration work can differ between objectives.

Each dataset contributes $4\times14\times3\times3=504$ runs.
The primary aggregate includes every retained fold/seed cell, giving
three seeds in each fold. The cap is the same numerical allowance on
both datasets, but their metrics concern different label vocabularies
and their starting difficulty need not be equal. Calibration remains
additional work outside this allowance.

Table~\ref{tab:cap-training-exposure} translates the admitted updates
into completed epochs. Here an epoch is one complete pass through the
selected training subset: ten crops for MSD and ADNI, or 19 scenes for
GeoFWI3D. Training visits each example once in a shuffled order before
starting the next pass, with batch size one. For $U$ training updates
and $N$ selected examples, the number of completed epochs and the
updates in the unfinished next epoch are
\begin{equation}
 E=\left\lfloor\frac{U}{N}\right\rfloor,
 \qquad R=U-NE,\qquad 0\leq R<N.
\label{eq:completed-training-epochs}
\end{equation}
Thus, 89 medical training updates are reported as eight completed
epochs plus nine further updates. Calibration evaluations do not count
as training epochs. These counts describe the training allowed by the
FLOP cap; an epoch limit is not imposed in addition to it.


% TABLE PROVENANCE: reports/thesis_noaug_update_20261021/results.json
% (MSD_capped and ADNI_capped) and
% reports/thesis_geofwi_update_20261022/results.json; 9 baseline and
% 117 constrained records per dataset/model. Epochs use integer division
% of training updates by 10 medical crops or 19 geological scenes.

Validation is recorded at predefined compute milestones and at the last
admissible state. Both best and final summaries are stored, but the
principal analysis uses the final state and has no early stopping.
Identifying the highest-mean objective after aggregation is a descriptive
selection across objectives, separate from checkpoint selection within
a run.

Saved final reports link configurations, manifests, charge ledgers and
run scores. The audit in Appendix~\ref{sec:evidence-register} distinguishes
recalculated aggregates from checks of the underlying archive payloads.
Disposed tensors and prediction volumes cannot be recovered from a
summary table; further inference requires a retained matching state
or a rerun.

\subsection{Geological Budget and Schedules}\label{sec:geological-budget-details}

The geological cap follows the same design principle as the medical
comparison: start from a short, practically useful update allowance for
nnU-Net, the most computationally expensive of the four implementations
per update, and give every architecture the corresponding FLOP allowance.
This choice aims to allow all models to learn within a restricted budget
despite their different computational costs. Models with cheaper updates
can consequently make more passes through the training subset. The
reference is update cost rather than parameter count; these two quantities
do not rank the architectures in the same order.

For the $96\times96\times96$ geological inputs, the reference allowance
was 152 nnU-Net updates, corresponding to eight complete passes through
19 training scenes. Multiplying this count by the original Dice-only
nnU-Net profile of $1{,}198.540306205$ GFLOPs per update established the
numerical cap. The final comparison retains it and charges both training
cost $C_{\mathrm{train}}$ and calibration cost $C_{\mathrm{cal}}$ within it:
% CAP PROVENANCE: scripts/cross_dataset_campaign/prepare_full.py sets
% SEISMIC_CAP from 152 times the original Dice-only nnU-Net update profile.
% The final native Dice+CE recipe retains that numerical cap.
\begin{equation}
 C_{\mathrm{train}}+C_{\mathrm{cal}}
 \leq 182{,}178.12654316\;\mathrm{GFLOPs}.
\label{eq:geological-compute-cap}
\end{equation}
This differs both in magnitude and accounting from the medical
optimization-only cap. Evaluation, diagnostic preflight and recipe
development remain outside this allowance. It is therefore not an
end-to-end compute budget, and absolute medical/geological scores do
not represent equal-compute comparisons across datasets.

Calibration uses all 19 training scenes at a fixed network state, without
augmentation or validation inputs. The sequential gradient-ratio procedure
uses median-target/95th-percentile limits of 3.6/4.5 for nnU-Net, 0.4/2.0
for UNETR and 0.1/0.5 for SwinUNETR and the compact U-Net. The nnU-Net and
UNETR strengths were selected following exploratory comparisons; they are
not shared coefficients across architectures. For combinations, sequential
calibration can also change a term's coefficient when another term is added.
The comparisons therefore evaluate complete calibrated configurations.

The unconstrained prefix and auxiliary ramp each last nine updates
for nnU-Net and 50 for the other models. SDF uses the early shared
ramp for SwinUNETR and UNETR in $S$, $B+S$ and $B+E+S$, and for
the compact U-Net in $B+S$. All other SDF configurations use the
frozen delayed policy, whose onset is placed near five ninths of the
nominal cap on the training trajectory, with a separate ramp of the
same length. The source schedule determines the exact integer update.
Learning-rate horizons stay fixed rather than being compressed to the
last affordable update.

The final native Dice-plus-cross-entropy nnU-Net recipe has a slightly
higher measured update cost than the original Dice-only profile. It
therefore fits 151 baseline updates within the retained cap: seven
completed epochs and 18 updates into the next, rather than the eight
complete epochs used to define the reference allowance. Reserving
calibration and charging auxiliary arithmetic further reduces the
training updates available to constrained configurations.
Table~\ref{tab:cap-training-exposure} reports the actual update counts,
completed epochs and remaining updates for every architecture under
both dataset-specific caps. Equal FLOPs specify the resource constraint;
the architectures can still be at different stages of learning.
The final admissible checkpoint supplies the reported score;
validation-best checkpoints and intermediate screening scores do not
replace it. Results appear in Section~\ref{sec:geo-results}.

\subsection{Model Settings and Training Exposure}\label{sec:implementation-tables}
These tables give the configured network sizes, activation policies and
training exposure underlying the main-text comparisons. Stochastic
gradient descent (SGD) updates parameters from batch gradients; momentum
also uses previous update directions. AdamW adjusts step sizes using
running gradient statistics and applies weight decay, a shrinkage of the
weights. A schedule horizon is the update count over which the chosen
learning-rate trajectory is defined.

\begin{table}[htbp]
\centering\small\setstretch{1.08}
\begin{tabular}{@{}P{2.5cm}rP{3.0cm}rr@{}}
\hline
Model & Parameters & Optimizer & Learning rate & Schedule horizon \\
\hline
nnU-Net & 16,548,076 & SGD, momentum 0.95 & $10^{-2}$ & 174 \\
SwinUNETR & 15,703,029 & AdamW & $10^{-4}$ & 610 \\
UNETR & 92,685,555 & AdamW & $10^{-4}$ & 477 \\
3D U-Net & 1,438,819 & AdamW & $3\times10^{-4}$ & 1,095 \\
\hline
\end{tabular}
\caption{MSD model parameter counts and fixed learning-rate horizons in optimizer
updates. nnU-Net uses Nesterov stochastic gradient descent (SGD) and weight decay $3\times10^{-5}$;
AdamW uses weight decay $10^{-5}$. A horizon specifies the trajectory,
not the number of updates available under the shorter FLOP cap. ADNI
uses binary output heads.}
\label{tab:main-model-settings}
\end{table}
\FloatBarrier
\begin{table}[htbp]
\centering
\small
\setstretch{1.05}
\renewcommand{\arraystretch}{1.2}
\begin{tabular}{@{}P{3.0cm}P{2.0cm}P{2.2cm}P{2.7cm}P{3.8cm}@{}}
\hline
Architecture & $S$ & $B+S$ & $B+E+S$ & Any included $S+D$ combination \\
\hline
nnU-Net & Delayed & Delayed & Delayed & Delayed \\
SwinUNETR & Shared & Shared & Shared & Delayed \\
UNETR & Shared & Shared & Shared & Delayed \\
3D U-Net & Delayed & Shared & Delayed & Delayed \\
\hline
\end{tabular}
\caption{Signed-distance activation policies frozen for the main
replication. The first three objective columns exclude depth. Shared
uses the early ramp; delayed starts near $5/9$ of the optimization
FLOP ceiling. Every entry also includes the supervised baseline.}
\label{tab:sdf-activation-policies}
\end{table}
\FloatBarrier
\begin{table}[htbp]
\centering
\small
\setstretch{1.08}
\setlength{\tabcolsep}{3.5pt}
\renewcommand{\arraystretch}{1.15}
\begin{tabular}{@{}llrrrrrr@{}}
\hline
 & & \multicolumn{3}{c}{Baseline} & \multicolumn{3}{c}{With constraints} \\
\cline{3-5}\cline{6-8}
Dataset & Model & Updates & Epochs & Extra & Updates & Epochs & Extra \\
\hline
\multicolumn{8}{@{}l}{Medical cap: $31{,}976.45066709$ GFLOPs; optimization only} \\
MSD & nnU-Net & 90 & 9 & 0 & 89 & 8 & 9 \\
MSD & SwinUNETR & 206 & 20 & 6 & 206 & 20 & 6 \\
MSD & UNETR & 245 & 24 & 5 & 245 & 24 & 5 \\
MSD & 3D U-Net & 563 & 56 & 3 & 562--563 & 56 & 2--3 \\
\hline
ADNI & nnU-Net & 90 & 9 & 0 & 89 / 90 & 8 / 9 & 9 / 0 \\
ADNI & SwinUNETR & 206 & 20 & 6 & 206 & 20 & 6 \\
ADNI & UNETR & 245 & 24 & 5 & 245 & 24 & 5 \\
ADNI & 3D U-Net & 563 & 56 & 3 & 562--563 & 56 & 2--3 \\
\hline
\multicolumn{8}{@{}l}{Geological cap: $182{,}178.12654316$ GFLOPs; training and calibration} \\
GeoFWI3D & nnU-Net & 151 & 7 & 18 & 145 & 7 & 12 \\
GeoFWI3D & SwinUNETR & 362 & 19 & 1 & 356 & 18 & 14 \\
GeoFWI3D & UNETR & 414 & 21 & 15 & 407--408 & 21 & 8--9 \\
GeoFWI3D & 3D U-Net & 952 & 50 & 2 & 943--945 & 49 & 12--14 \\
\hline
\end{tabular}
\caption{Training exposure under the two FLOP caps, calculated from the
completed three-fold, three-seed runs. Epochs are whole completed passes
through the selected training subset; Extra gives updates in the
unfinished next pass. Ranges cover the thirteen constrained objectives.
For ADNI nnU-Net, slash-separated entries correspond in order: 89 updates
mean eight epochs plus nine updates, while 90 mean nine epochs with no
remainder. Calibration is excluded from these epoch counts, even when
its computation is charged within the geological cap.}
\label{tab:cap-training-exposure}
\end{table}
\FloatBarrier


\subsection{Signed-Distance Field Construction}
\label{sec:sdf-construction-details}

For $M_{i,c}=\{v\in\Omega_i:y_i(v)=c\}$, define
$\operatorname{dist}_{\mathrm{vox}}(v,U)=\min_{u\in U}\|v-u\|_2$
on the unit voxel grid. When the class and its complement are nonempty,
the signed field is
\begin{equation}
\widetilde\varphi_{i,c}(v)=
\begin{cases}
-\bigl(\operatorname{dist}_{\mathrm{vox}}
       (v,\Omega_i\setminus M_{i,c})-\tfrac12\bigr),
   &v\in M_{i,c},\\[0.4em]
\operatorname{dist}_{\mathrm{vox}}(v,M_{i,c})-\tfrac12,
   &v\notin M_{i,c}.
\end{cases}
\label{eq:sdf-half-voxel-field}
\end{equation}
It is negative inside the class and positive outside. The half-voxel
shift places adjacent inside and outside centres on opposite sides of
the interface. This defines a discrete distance convention; it is not
an exact continuous-surface distance for every boundary orientation.
Clipping replaces values outside $[-\tau,\tau]$ by the corresponding
endpoint, and division by $\tau$ normalizes their scale:
\begin{equation}
\varphi_{i,c}(v)
=\frac{\operatorname{clip}
  (\widetilde\varphi_{i,c}(v),-\tau,\tau)}{\tau},
\qquad \tau=8.
\label{eq:sdf-clipped-normalized-field}
\end{equation}
so geometric weights are bounded. An absent class is assigned $+1$
throughout the grid and a full-grid class $-1$; both are retained in
the class reduction. The field uses the stored grid without adding an
exterior background frame.


\subsection{Spatial Supervision on Actual Dataset Volumes}
\label{sec:dataset-constraint-examples}

Figures~\ref{fig:msd-constraint-example}--\ref{fig:geofwi3d-constraint-example}
use the same cases as the dataset illustrations: \texttt{hippocampus\_001},
\texttt{hippocampus\_adni\_10007\_L} and \texttt{model\_0001}.
Their native grids contain $35\times51\times35$, $32\times48\times40$
and $96\times96\times96$ voxels, respectively. All morphology, distance
fields, depth layers and counts are computed in three dimensions before
rendering. The examples illustrate the definitions on unaugmented native
grids; training applies the constraints to the grid specified by each
recipe, which can change the available surrounding region and its weights.

Panel (a) shows the intact three-dimensional annotation. Panels (b) and (c)
show two-step, six-neighbour bands and the in-plane subset of their selected
voxel pairs. Bands and edges use the whole foreground, including the union
of anterior and posterior hippocampus in MSD. Blue pair segments connect
equal foreground-membership labels; orange segments cross between foreground
and background. The anatomical division within MSD is therefore absent from
this pair target. The geological example retains zero-padded morphology at
the cube faces, as in three of the four conventional model recipes. UNETR's
frame exclusion is specified in Appendix~\ref{sec:geo-boundaries}.

The band view is the source-axis-2 slice containing the most foreground
voxels, with the lowest index breaking ties. The yellow square identifies
the $16\times16$ area enlarged in the pair view. Its centre is placed near
the in-plane foreground boundary closest to the slice's foreground centroid,
away from the crop faces where possible. Pairs extending out of the displayed
plane still enter the three-dimensional objective but are not drawn.

Panel (d) renders the signed-distance field of anterior hippocampus for MSD
and of the single foreground class for ADNI and GeoFWI3D. Distances use the
half-voxel convention, clipping at eight voxels and division by eight. The
rendering retains the class and exterior field values up to four voxel units;
a quadrant is cut away to expose the interior. These display choices do not
alter the computed field. The intact class mask in (a) provides its context.

Panels (e) and (f) explain depth balancing. The bars in (e) show each layer's
share of the surrounding region's voxels; the line shows its share of the
loss weight. The latter is equal across layers with at least 64 voxels and
smaller for sparser layers. Panel (f) shows five occupied layers sampled
uniformly across the ordered layer indices, including both extremes. Each
retains its full three-dimensional shape and is translated along the plane
normal to separate it from its neighbours. Colours use the voxel weights
computed from \emph{all} layers, divided by their mean over the full region;
they do not depict prediction errors. The common colour scale is logarithmic.

MSD uses its verified physical coordinate map, with two-millimetre layers
and an eight-millimetre surrounding margin. ADNI and GeoFWI3D use two-voxel
layers and an eight-voxel margin in source-array coordinates. The reference
plane follows the dataset-specific construction described in
Chapter~\ref{sec:methodology}. None of these examples requires the geological
fallback for a nearly parallel principal axis. Volumes are rotated for
presentation without smoothing or resampling their voxel faces; rendered
sizes are fitted independently and do not imply a common physical scale.

\begin{figure}[p]
\centering
\includegraphics[width=\linewidth]{figures/compiled/msd-constraint-example.pdf}
\caption{Spatial supervision on an MSD volume. The intact annotation in
(a) distinguishes anterior (blue) and posterior (orange) hippocampus.
The union bands in (b) and enlarged neighbour pairs in (c) omit this
internal division, whereas (d) shows a distance field specific to the
anterior class. The full-volume layer distribution in (e) explains the
weights of the sampled three-dimensional layers in (f): one sparse layer
receives less total weight, while the remaining layers have equal shares.
The yellow square in (b) locates the enlarged pair view.}
\label{fig:msd-constraint-example}
\end{figure}

\begin{figure}[p]
\centering
\includegraphics[width=\linewidth]{figures/compiled/adni-constraint-example.pdf}
\caption{Spatial supervision on an ADNI hippocampus volume. Panels follow
the same arrangement as Figure~\ref{fig:msd-constraint-example}, with one
foreground class. The three-dimensional signed-distance cutaway shows its
interior and surrounding background. Every occupied depth layer contains
at least 64 voxels in this example, so the layer-level loss shares in (e)
are equal despite different voxel counts. The smaller layers therefore
have larger weights per voxel, visible in (f).}
\label{fig:adni-constraint-example}
\end{figure}

\begin{figure}[p]
\centering
\includegraphics[width=\linewidth]{figures/compiled/geofwi3d-constraint-example.pdf}
\caption{Spatial supervision on a GeoFWI3D salt volume. The three-dimensional
annotation in (a) and distance cutaway in (d) retain the native salt geometry.
Bands and pairs in (b)--(c) follow its observed interfaces under zero-padded
morphology. All occupied depth layers exceed the sparse-layer threshold,
giving equal layer-level loss shares in (e). The separated layers in (f)
use the same annotation-derived weighting rule as the medical examples,
in voxel units; their ordering does not represent geological age.}
\label{fig:geofwi3d-constraint-example}
\end{figure}
\FloatBarrier

\subsection{Conventional Model and Optimizer Settings}
\label{sec:conventional-implementation}

Here a layer's width is its number of feature channels; embedding width is
the length of each token vector. Pooling reduces spatial resolution, and
transposed convolutions perform learned upsampling. Concatenated skips join
encoder and decoder channels, while residual blocks add a block's input to
its transformed output. Instance normalization rescales each case's feature
channels; leaky rectified linear activations retain a small slope for
negative inputs. Multiple attention heads apply separate attention
calculations within a layer.

The compact model uses MONAI's \texttt{BasicUNet}, with encoder widths
$16,16,32,64,128$ and final decoder width 16. It uses pooling, transposed
convolutions and concatenated skips, without residual blocks or deep
supervision. The fixed nnU-Net uses five stages of widths
$32,64,128,256,320$, three-dimensional kernels, instance normalization and
leaky rectified linear activations. UNETR has 12 blocks, 768-dimensional
embeddings, 12 attention heads and $16^3$ patches, giving 64 tokens per
crop; its residual convolutional decoder has base width 16. SwinUNETR
combines a hierarchical shifted-window encoder with a convolutional decoder.


A warmup gradually increases the learning rate at the start of training;
decay lowers it later. nnU-Net uses polynomial learning-rate decay with
exponent 0.9 and no warmup. AdamW models warm up over 5\% of their fixed horizon and then
decay cosinusoidally toward 1\% of the initial rate. Every objective
for a model follows the same trajectory. When an extra loss reduces
the number of affordable updates, the learning-rate horizon remains
unchanged. Parameters and optimizer states use float32, eligible CUDA
operations use bfloat16 autocast, and constraints retain their stable
float32 calculations.

\subsection{Pretrained Adaptation Settings and Execution Records}
\label{sec:adaptation-implementation}

Section~\ref{sec:sam3-methods} defines the SAM3 support data, input
preparation and binary prediction rule; Sections~\ref{sec:medsam-methods}
and \ref{sec:medsam-adni-methods} give the corresponding MedSAM3 protocols.
This section records their exact optimizer, calibration and execution
settings. Gradient-norm clipping rescales an excessively large gradient
vector to a maximum allowed length. Accumulation combines gradients from
sequentially processed slice pairs before one parameter update. A slab here
consists of an anchor plane and its adjacent companion.

\paragraph{SAM3 on GeoFWI3D.}\label{sec:sam3-methods-details}
Folds 0, 1 and 3 are crossed with seeds 0, 1 and 2 and fourteen
objectives, giving 126 objective runs. Each update accumulates two
adjacent-slice pairs. The 25 updates therefore give 100 slice and
single-class-query presentations, repeatedly drawn from five support
volumes. AdamW uses weight decay 0.01 and gradient-norm clipping at
one. Its retained 200-update learning-rate horizon has a 20-update
warmup and peak/final rates of $5\times10^{-5}$ and $5\times10^{-6}$.
Only the first 25 steps are executed.

Calibration uses four anchor pairs per support volume, giving twenty
probes. Bands are calibrated at initialization for each fold/seed cell.
Other selected terms are calibrated sequentially after five prefix
updates, in the order $B\to E\to S\to D$, retaining the band
coefficient when present. Each additional term targets a median
gradient ratio of 0.1 with a 95th-percentile cap of 0.5. Terms awaiting
calibration have zero weight during the prefix; their fixed calibrated
coefficients then follow the global ramp $\min\{1,t/20\}$. A term
entering at update 6 consequently receives ramp factor $6/20$.

The sampled salt losses use full-volume supports and counts. Bands
retain two-step zero-padded morphology, including crop faces.
The half-voxel signed field uses unit-grid distances clipped at eight
voxels and divided by eight; depth uses source-array geometry.
The sampled reductions are in
Appendix~\ref{sec:sampled-pretrained-supervision}.

Each arm restores the same base and initial adapters and uses its
matched sampling schedule. Evaluation after update 25 covers all five
support and 96 held-out volumes. Paired gains use only the SAM3
baseline from the same fold and seed.

The completed inventory contains 76 runs on Colab A100-SXM4-80GB graphics
processing units (GPUs) and 50 on A100 80GB PCIe GPUs allocated as 4g.40gb
Multi-Instance GPU (MIG) partitions, allocated subdivisions of a physical
GPU. All retain the same scientific source identity. Archived hashes,
digital fingerprints of file or array contents, verify unchanged base
parameters and matching initial adapters, support arrays and sampling
schedules within each comparison. Hardware allocation was not fully crossed
with the objectives, so timing records retain their execution platform.
Recovery attempts are distinguished from the final 25-update trajectory.

\paragraph{MedSAM3 on MSD.}\label{sec:medsam-methods-details}
Folds 0, 1 and 3 are crossed with seeds 17, 83 and 191. Fourteen
objectives give 126 completed objective runs across the
nine fold/seed cells, with repeated matched controls retained in the
execution records. The 916 trainable adapter tensors (multidimensional parameter arrays)
start from the
released medical factors. Each 200-update run uses two adjacent-slice pairs per update,
giving 800 slice and 1,600 foreground-query presentations. These
repetitions do not increase the five-crop annotation budget.

AdamW uses 20 warmup updates, peak/final rates of $5\times10^{-5}$
and $5\times10^{-6}$, weight decay 0.01 and gradient-norm clipping
at one. Every selected auxiliary term ramps as $\min\{1,t/20\}$
from the first update, without the delayed SDF policy of the conventional
models. Support volumes and anchors are sampled uniformly. Bands,
SDF and depth use the anchor plane; edges also use its companion.
The full-volume denominators and sampling factors are specified in
Appendix~\ref{sec:sampled-pretrained-supervision}.

Calibration uses four anchors per support volume, giving twenty
slabs. Gradients are measured with respect to the LoRA parameters,
retaining vector interactions in the sequential order
$B\to E\to S\to D$. Equation~\ref{eq:constraint-calibrated-coefficient}
uses $(\eta_a,\kappa_a)=(0.1,0.5)$. Coefficients depend on the selected
combination and use support data only. Within each fold they are shared
across optimization seeds according to the saved calibration-source
records. The $B$ and $B+E$ variants retain their original band and edge
calibration and matched baseline; the other eleven constrained variants
use their separately audited execution batch. Coefficients and controls
are not substituted between these groups.

Bands and edges average the background, anterior and posterior masks.
SDF uses the foreground fields with half-voxel shifts, clipping at eight
units and division by eight on MSD's verified 1 mm grid. Depth uses
physical geometry. Native-grid inference and final-update evaluation
follow Section~\ref{sec:medsam-methods}.

\paragraph{MedSAM3 on ADNI.}\label{sec:medsam-adni-methods-details}
Folds 0, 1 and 3 and seeds 0, 1 and 2 give 126 runs across fourteen
objectives. Support selection follows the saved subject manifest and
alternates left/right crops; the overlap between support sets is
specified in Section~\ref{sec:medsam-adni-methods}. The medical
initialization and optimizer settings are the same as for MSD.
With a single foreground query, the 200 updates give 800 slice and
800 query presentations from two support crops.

Four anchors per support volume give eight calibration slabs.
The sequential gradient calibration uses
$(\eta_a,\kappa_a)=(0.1,0.5)$ with combination-specific coefficients
computed from support data only. Schedule hashes verify matched
sampling across arms. Bands and edges use the single hippocampus
mask and its complement, without a second background-mask mean.
The half-voxel SDF is clipped at eight unit array voxels and divided
by eight. Depth uses the source-axis geometric reference.
Evaluation uses the final 200-update adapter on both support crops
and all 54 held-out crops.

In both medical studies, volume scores and mean per-crop confusion counts
are averaged within runs, across seeds within a fold and equally across
folds. No early stopping or validation-dependent checkpoint selection is
used. Training-loop timers exclude initialization, calibration, validation
and inference, and provide no peak-memory measurement. The adaptation
allowance excludes pretraining. A common seed alone does not guarantee
bitwise determinism, or numerically identical stored outputs, in CUDA
execution. Execution and calibration records are indexed in
Appendix~\ref{sec:evidence-register}.

\clearpage
\section{Supplementary Experimental Results}
\label{sec:extended-results}

This appendix retains the full numerical comparisons behind the main
results figures. Scores are percentages, $\Delta$ is a paired change in
percentage points, and $n_+$ counts positive changes among the nine
fold/seed pairs. These pairs share data and are not nine independent
datasets. Standard deviations (SDs) describe the three seed-averaged fold
means. Training exposure is reported in
Table~\ref{tab:cap-training-exposure}; individual run records are indexed
in Appendix~\ref{sec:evidence-register}.

\subsection{Complete Dice Comparisons}\label{sec:complete-result-tables}
The following tables preserve all evaluated objectives and their matched
controls. Main-text figures summarize selected comparisons; these tables
retain the complete numerical record, including negative results.

\begin{table}[htbp]
\centering\small\setstretch{1.05}
\setlength{\tabcolsep}{5pt}
\renewcommand{\arraystretch}{1.15}
\begin{tabular}{@{}lrrrr@{}}
\hline
 & \multicolumn{4}{c}{Validation Dice (\%)} \\
Training variant & nnU-Net & SwinUNETR & UNETR & 3D U-Net \\
\hline
Baseline & $72.08\pm1.45$ & $55.54\pm0.57$ & $54.91\pm2.30$ & $77.60\pm2.59$ \\
Baseline $+B$ & $73.43\pm1.29$ & $57.13\pm1.05$ & $55.66\pm1.63$ & $78.32\pm2.46$ \\
Baseline $+S$ & $72.01\pm1.58$ & $55.91\pm0.15$ & $54.57\pm3.05$ & $77.83\pm2.68$ \\
Baseline $+B+S$ & $73.14\pm2.05$ & $56.98\pm1.13$ & $54.97\pm3.28$ & $77.91\pm1.10$ \\
Baseline $+B+E$ & $72.84\pm1.41$ & $56.99\pm0.96$ & $\mathbf{56.07\pm2.19}$ & $78.14\pm1.58$ \\
Baseline $+B+E+S$ & $71.79\pm3.67$ & $56.93\pm1.05$ & $54.70\pm3.32$ & $78.56\pm1.94$ \\
Baseline $+D$ & $73.92\pm1.08$ & $56.55\pm0.35$ & $55.59\pm2.66$ & $79.40\pm3.31$ \\
Baseline $+B+D$ & $75.23\pm2.65$ & $57.28\pm0.71$ & $55.69\pm1.60$ & $80.40\pm2.37$ \\
Baseline $+E+D$ & $75.21\pm2.37$ & $56.76\pm0.66$ & $55.68\pm2.45$ & $79.59\pm2.97$ \\
Baseline $+S+D$ & $76.26\pm2.20$ & $56.50\pm0.47$ & $55.12\pm2.94$ & $79.77\pm2.95$ \\
Baseline $+B+E+D$ & $70.72\pm3.39$ & $57.33\pm1.17$ & $55.85\pm1.99$ & $80.50\pm2.03$ \\
Baseline $+B+S+D$ & $77.33\pm2.19$ & $57.58\pm0.95$ & $55.31\pm2.78$ & $80.73\pm2.31$ \\
Baseline $+E+S+D$ & $76.85\pm2.96$ & $56.76\pm0.53$ & $54.67\pm2.96$ & $80.28\pm2.79$ \\
Baseline $+B+E+S+D$ & $\mathbf{78.22\pm2.42}$ & $\mathbf{57.61\pm0.83}$ & $55.55\pm3.35$ & $\mathbf{80.73\pm2.24}$ \\
\hline
\end{tabular}
\caption{GeoFWI3D final-budget salt Dice (\%). Each condition contains
three folds and three seeds. Dispersion is the sample SD of the three
seed-averaged fold means. Bold marks the highest mean within each model.
The $182{,}178.13$ GFLOP cap includes calibration. UNETR uses
Dice plus cross-entropy with horizontal rotations and reflections during
training.}
\label{tab:geo-matrix}
\end{table}
\FloatBarrier
\begin{table}[htbp]
\centering\footnotesize\setstretch{1.05}
\setlength{\tabcolsep}{5pt}
\begin{tabular}{@{}llrrrl@{}}
\hline
Model & Variant & $\Delta$ & SD($\Delta_f$) & Positive & Fold 0 / 1 / 2 gains \\
\hline
nnU-Net & Baseline $+B+E+S+D$ & +6.14 & 1.00 & 9/9 & +6.71 / +6.71 / +4.99 \\
SwinUNETR & Baseline $+B+E+S+D$ & +2.07 & 1.31 & 9/9 & +3.46 / +0.86 / +1.88 \\
UNETR & Baseline $+B+E$ & +1.16 & 0.33 & 8/9 & +1.16 / +1.49 / +0.83 \\
3D U-Net & Baseline $+B+E+S+D$ & +3.14 & 1.84 & 9/9 & +4.44 / +1.03 / +3.93 \\
\hline
\end{tabular}
\caption{Paired geological gains in percentage points for each model's
highest-mean objective. All three fold-mean gains are positive; UNETR
has one non-improving seed cell. Selection of the highest means is
descriptive and uses the same validation results being reported.}
\label{tab:geo-best}
\end{table}
\FloatBarrier
\begin{table}[htbp]
\centering\small\setstretch{1.05}
\setlength{\tabcolsep}{5pt}
\renewcommand{\arraystretch}{1.15}
\begin{tabular}{@{}lrrrr@{}}
\hline
 & \multicolumn{4}{c}{Validation Dice (\%)} \\
Training variant & nnU-Net & SwinUNETR & UNETR & 3D U-Net \\
\hline
Baseline & $62.52\pm7.37$ & $54.04\pm1.53$ & $54.71\pm0.50$ & $73.83\pm1.08$ \\
Baseline $+B$ & $70.07\pm3.13$ & $57.89\pm1.28$ & $58.60\pm0.91$ & $78.66\pm1.07$ \\
Baseline $+S$ & $60.48\pm7.98$ & $54.30\pm1.38$ & $55.49\pm0.11$ & $73.66\pm1.54$ \\
Baseline $+B+S$ & $70.12\pm3.23$ & $59.01\pm1.58$ & $59.86\pm0.70$ & $78.68\pm0.86$ \\
Baseline $+B+E$ & $68.97\pm3.76$ & $58.03\pm1.31$ & $58.85\pm0.83$ & $79.54\pm0.78$ \\
Baseline $+B+E+S$ & $69.00\pm3.83$ & $59.20\pm1.52$ & $\mathbf{60.07\pm0.33}$ & $79.36\pm0.64$ \\
Baseline $+D$ & $64.13\pm6.50$ & $62.28\pm1.29$ & $56.47\pm1.12$ & $76.13\pm1.18$ \\
Baseline $+B+D$ & $71.90\pm2.63$ & $65.07\pm1.05$ & $58.42\pm1.31$ & $78.85\pm0.70$ \\
Baseline $+E+D$ & $62.98\pm7.27$ & $63.00\pm1.37$ & $56.92\pm0.96$ & $77.02\pm1.05$ \\
Baseline $+S+D$ & $65.52\pm6.04$ & $62.48\pm1.39$ & $57.08\pm0.77$ & $76.19\pm0.91$ \\
Baseline $+B+E+D$ & $71.61\pm2.99$ & $65.59\pm0.98$ & $58.52\pm1.14$ & $\mathbf{80.00\pm0.64}$ \\
Baseline $+B+S+D$ & $\mathbf{72.67\pm2.55}$ & $65.70\pm0.96$ & $59.41\pm0.98$ & $78.78\pm1.15$ \\
Baseline $+E+S+D$ & $64.91\pm6.68$ & $63.16\pm1.19$ & $57.42\pm0.83$ & $77.16\pm1.32$ \\
Baseline $+B+E+S+D$ & $72.50\pm2.89$ & $\mathbf{66.00\pm1.02}$ & $59.60\pm0.89$ & $79.89\pm0.87$ \\
\hline
\end{tabular}
\caption{MSD final optimization-budget Dice (\%). Each entry averages three seeds per fold and three folds equally; dispersion is across fold means. Bold marks the highest mean within each model, selected descriptively.}
\label{tab:msd-noaug-matrix}
\end{table}
\FloatBarrier
\begin{table}[htbp]
\centering\footnotesize\setstretch{1.05}
\setlength{\tabcolsep}{5pt}
\renewcommand{\arraystretch}{1.15}
\begin{tabular}{@{}llrrrl@{}}
\hline
Model & Variant & $\Delta$ & SD($\Delta_f$) & Positive & Fold 0 / 1 / 2 gains \\
\hline
nnU-Net & Baseline $+B+S+D$ & +10.15 & 4.87 & 9/9 & +15.66 / +6.41 / +8.39 \\
SwinUNETR & Baseline $+B+E+S+D$ & +11.97 & 0.72 & 9/9 & +11.63 / +12.79 / +11.47 \\
UNETR & Baseline $+B+E+S$ & +5.35 & 0.81 & 9/9 & +4.66 / +6.24 / +5.16 \\
3D U-Net & Baseline $+B+E+D$ & +6.17 & 0.52 & 9/9 & +5.92 / +6.77 / +5.81 \\
\hline
\end{tabular}
\caption{Paired changes for the highest-mean MSD configuration of each model. Fold gains average the three paired seed differences. The nine cells share data and are not independent test datasets.}
\label{tab:msd-noaug-best}
\end{table}
\FloatBarrier
\begin{table}[htbp]
\centering\small\setstretch{1.05}
\setlength{\tabcolsep}{5pt}
\renewcommand{\arraystretch}{1.15}
\begin{tabular}{@{}lrrrr@{}}
\hline
 & \multicolumn{4}{c}{Validation Dice (\%)} \\
Training variant & nnU-Net & SwinUNETR & UNETR & 3D U-Net \\
\hline
Baseline & $75.66\pm1.49$ & $68.72\pm2.86$ & $56.27\pm2.41$ & $69.64\pm3.13$ \\
Baseline $+B$ & $76.39\pm1.61$ & $72.06\pm2.42$ & $59.26\pm1.89$ & $73.45\pm2.32$ \\
Baseline $+S$ & $75.67\pm1.47$ & $68.65\pm2.66$ & $55.92\pm2.27$ & $69.40\pm3.18$ \\
Baseline $+B+S$ & $76.40\pm1.62$ & $71.99\pm2.56$ & $59.51\pm1.82$ & $74.31\pm2.55$ \\
Baseline $+B+E$ & $76.25\pm1.65$ & $72.40\pm2.34$ & $59.76\pm1.95$ & $74.21\pm3.87$ \\
Baseline $+B+E+S$ & $76.26\pm1.68$ & $72.34\pm2.51$ & $60.18\pm1.86$ & $74.40\pm3.55$ \\
Baseline $+D$ & $75.97\pm1.53$ & $70.07\pm2.56$ & $57.77\pm1.96$ & $71.61\pm2.38$ \\
Baseline $+B+D$ & $76.62\pm1.52$ & $72.63\pm2.34$ & $59.97\pm1.78$ & $74.15\pm2.91$ \\
Baseline $+E+D$ & $76.14\pm1.46$ & $70.50\pm2.66$ & $58.12\pm2.24$ & $71.93\pm2.76$ \\
Baseline $+S+D$ & $76.09\pm1.49$ & $70.03\pm2.66$ & $57.79\pm1.87$ & $71.62\pm2.38$ \\
Baseline $+B+E+D$ & $76.53\pm1.58$ & $\mathbf{73.03\pm2.27}$ & $60.36\pm1.59$ & $\mathbf{75.18\pm2.64}$ \\
Baseline $+B+S+D$ & $\mathbf{76.69\pm1.44}$ & $72.52\pm2.40$ & $59.79\pm1.86$ & $74.36\pm3.08$ \\
Baseline $+E+S+D$ & $76.19\pm1.55$ & $70.48\pm2.66$ & $58.34\pm2.20$ & $72.39\pm3.12$ \\
Baseline $+B+E+S+D$ & $76.54\pm1.48$ & $72.94\pm2.40$ & $\mathbf{60.46\pm1.43}$ & $75.02\pm2.99$ \\
\hline
\end{tabular}
\caption{ADNI final optimization-budget Dice (\%). Each entry averages three seeds per fold and three folds equally; dispersion is across fold means. Bold marks the highest mean within each model, selected descriptively.}
\label{tab:adni-noaug-matrix}
\end{table}
\FloatBarrier
\begin{table}[htbp]
\centering\footnotesize\setstretch{1.05}
\setlength{\tabcolsep}{5pt}
\renewcommand{\arraystretch}{1.15}
\begin{tabular}{@{}llrrrl@{}}
\hline
Model & Variant & $\Delta$ & SD($\Delta_f$) & Positive & Fold 0 / 1 / 2 gains \\
\hline
nnU-Net & Baseline $+B+S+D$ & +1.03 & 0.17 & 9/9 & +1.19 / +1.05 / +0.86 \\
SwinUNETR & Baseline $+B+E+D$ & +4.31 & 0.86 & 9/9 & +5.18 / +4.27 / +3.46 \\
UNETR & Baseline $+B+E+S+D$ & +4.20 & 0.98 & 9/9 & +5.33 / +3.70 / +3.56 \\
3D U-Net & Baseline $+B+E+D$ & +5.54 & 0.58 & 9/9 & +6.18 / +5.42 / +5.03 \\
\hline
\end{tabular}
\caption{Paired changes for the highest-mean ADNI configuration of each model. Fold gains average the three paired seed differences. The nine cells share data and are not independent test datasets.}
\label{tab:adni-noaug-best}
\end{table}
\FloatBarrier
\begin{table}[htbp]
\centering\small\setstretch{1.05}
\setlength{\tabcolsep}{6pt}
\renewcommand{\arraystretch}{1.12}
\begin{tabular}{@{}lrrrr@{}}
\hline
\shortstack{Training\\variant} & \shortstack{Validation\\Dice (\%)} & \shortstack{$\Delta$ Dice\\(pp)} & \shortstack{Positive\\pairs} & \shortstack{Training-set\\Dice (\%)} \\
\hline
Baseline & 69.31 $\pm$ 4.00 & 0.00 & --- & 74.97 $\pm$ 5.73 \\
Baseline $+B$ & 70.79 $\pm$ 2.38 & +1.48 $\pm$ 4.16 & 6/9 & 77.03 $\pm$ 7.11 \\
Baseline $+S$ & 69.10 $\pm$ 4.88 & -0.22 $\pm$ 2.70 & 5/9 & 75.31 $\pm$ 7.92 \\
Baseline $+B+S$ & 69.60 $\pm$ 1.54 & +0.29 $\pm$ 3.75 & 3/9 & 76.67 $\pm$ 5.77 \\
Baseline $+B+E$ & 69.44 $\pm$ 3.45 & +0.12 $\pm$ 0.84 & 3/9 & 76.12 $\pm$ 6.57 \\
Baseline $+B+E+S$ & 67.46 $\pm$ 1.26 & -1.85 $\pm$ 2.79 & 3/9 & 76.14 $\pm$ 5.22 \\
Baseline $+D$ & 72.29 $\pm$ 2.03 & +2.98 $\pm$ 2.05 & 8/9 & 77.55 $\pm$ 6.22 \\
Baseline $+B+D$ & 72.45 $\pm$ 0.70 & +3.13 $\pm$ 3.30 & 7/9 & 77.47 $\pm$ 5.22 \\
Baseline $+E+D$ & 69.18 $\pm$ 2.97 & -0.13 $\pm$ 1.79 & 5/9 & 76.14 $\pm$ 4.90 \\
Baseline $+S+D$ & 69.93 $\pm$ 1.29 & +0.61 $\pm$ 3.91 & 5/9 & 76.34 $\pm$ 6.34 \\
Baseline $+B+E+D$ & 72.76 $\pm$ 1.97 & +3.45 $\pm$ 2.43 & 8/9 & 77.41 $\pm$ 6.86 \\
Baseline $+B+S+D$ & 72.17 $\pm$ 0.89 & +2.85 $\pm$ 3.13 & 7/9 & 77.40 $\pm$ 5.58 \\
Baseline $+E+S+D$ & 69.70 $\pm$ 3.23 & +0.39 $\pm$ 4.18 & 5/9 & 75.52 $\pm$ 2.88 \\
Baseline $+B+E+S+D$ & 70.95 $\pm$ 1.44 & +1.64 $\pm$ 2.56 & 7/9 & 77.79 $\pm$ 5.06 \\
\hline
\end{tabular}
\caption{SAM3 on GeoFWI3D after 25 updates. Both Dice columns report
hard salt overlap in percent. $\Delta$ Dice is the paired validation
Dice change relative to the matched baseline. Dice entries contain mean
$\pm$ sample SD across three seed-averaged fold means; paired changes
contain mean $\pm$ SD of the three fold-mean differences. Positive counts
use the nine paired runs. Each row represents the full three-fold,
three-seed condition, verified against the final archives;
baseline changes are zero by definition.}
\label{tab:sam3-geo-results}
\end{table}
\FloatBarrier
\begin{table}[htbp]
\centering\footnotesize\setstretch{1.05}
\setlength{\tabcolsep}{4pt}
\renewcommand{\arraystretch}{1.15}
\begin{tabular}{@{}lrrrrr@{}}
\hline
\shortstack{Training\\variant} & \shortstack{Matched baseline\\validation Dice\\(\%)} & \shortstack{Validation\\Dice (\%)} & \shortstack{$\Delta$ Dice\\(pp)} & \shortstack{Positive\\pairs} & \shortstack{Training-set\\Dice (\%)} \\
\hline
Baseline $+B$ & 74.93 & $78.76\pm0.36$ & +3.83 & 8/9 & 83.82 \\
Baseline $+S$ & 76.34 & $76.96\pm2.55$ & +0.62 & 6/9 & 81.49 \\
Baseline $+B+S$ & 76.34 & $80.60\pm1.20$ & +4.26 & 9/9 & 85.29 \\
Baseline $+B+E$ & 74.93 & $80.12\pm1.21$ & +5.19 & 9/9 & 85.80 \\
Baseline $+B+E+S$ & 76.34 & $81.02\pm1.49$ & +4.68 & 9/9 & 86.44 \\
Baseline $+D$ & 76.34 & $78.90\pm1.45$ & +2.56 & 8/9 & 84.21 \\
Baseline $+B+D$ & 76.34 & $80.45\pm0.90$ & +4.10 & 9/9 & 85.58 \\
Baseline $+E+D$ & 76.34 & $80.74\pm1.50$ & +4.39 & 9/9 & 85.84 \\
Baseline $+S+D$ & 76.34 & $79.84\pm1.00$ & +3.49 & 9/9 & 84.53 \\
Baseline $+B+E+D$ & 76.34 & $80.95\pm1.53$ & +4.60 & 9/9 & 86.29 \\
Baseline $+B+S+D$ & 76.34 & $81.09\pm1.21$ & +4.75 & 9/9 & 86.01 \\
Baseline $+E+S+D$ & 76.34 & $81.28\pm1.42$ & +4.94 & 9/9 & 86.40 \\
Baseline $+B+E+S+D$ & 76.34 & $81.73\pm1.43$ & +5.38 & 9/9 & 86.88 \\
\hline
\end{tabular}
\caption{MSD MedSAM3 after 200 updates. All Dice columns report the
mean of hard anterior and posterior Dice in percent. The matched-control
and validation columns both use held-out crops; the training-set column
uses support crops. $\Delta$ Dice is the paired validation gain over the
control in the same row. The baseline and thirteen constrained variants
complete the fourteen-objective inventory. Each variant includes three
folds and three seeds; dispersion is the standard deviation of the
three seed-averaged fold means.}
\label{tab:medsam-msd-noaug}
\label{tab:medsam-msd-original-noaug}
\end{table}
\FloatBarrier
\begin{table}[htbp]
\centering\small\setstretch{1.05}
\setlength{\tabcolsep}{5pt}
\renewcommand{\arraystretch}{1.15}
\begin{tabular}{@{}lrrrr@{}}
\hline
\shortstack{Training\\variant} & \shortstack{Validation\\Dice (\%)} & \shortstack{$\Delta$ Dice\\(pp)} & \shortstack{Positive\\pairs} & \shortstack{Training-set\\Dice (\%)} \\
\hline
Baseline & $77.32\pm1.95$ & +0.00 & --- & 90.62 \\
Baseline $+B$ & $79.08\pm2.81$ & +1.75 & 8/9 & 91.78 \\
Baseline $+S$ & $78.22\pm2.03$ & +0.90 & 9/9 & 91.12 \\
Baseline $+B+S$ & $80.36\pm2.10$ & +3.03 & 9/9 & 92.48 \\
Baseline $+B+E$ & $79.24\pm2.33$ & +1.91 & 7/9 & 92.08 \\
Baseline $+B+E+S$ & $80.65\pm1.84$ & +3.33 & 9/9 & 93.79 \\
Baseline $+D$ & $79.09\pm2.64$ & +1.77 & 8/9 & 91.66 \\
Baseline $+B+D$ & $80.18\pm2.28$ & +2.85 & 9/9 & 92.61 \\
Baseline $+E+D$ & $79.69\pm2.79$ & +2.37 & 9/9 & 92.59 \\
Baseline $+S+D$ & $79.23\pm2.34$ & +1.91 & 9/9 & 92.00 \\
Baseline $+B+E+D$ & $79.84\pm2.48$ & +2.51 & 9/9 & 93.05 \\
Baseline $+B+S+D$ & $80.38\pm1.99$ & +3.06 & 9/9 & 92.73 \\
Baseline $+E+S+D$ & $80.85\pm2.05$ & +3.53 & 9/9 & 93.86 \\
Baseline $+B+E+S+D$ & $80.74\pm2.05$ & +3.42 & 9/9 & 93.74 \\
\hline
\end{tabular}
\caption{ADNI MedSAM3 after 200 updates. Both Dice columns report hard
hippocampus overlap in percent: validation uses held-out crops and
training-set evaluation uses support crops. $\Delta$ Dice is the paired
validation gain over the baseline. Dispersion is across three
seed-averaged fold means. Each fold uses two support crops; folds 0 and 3
share the same support set.}
\label{tab:medsam-adni-noaug}
\end{table}
\FloatBarrier
\subsection{Additional Objective-Level Comparisons}

\label{sec:secondary-objective-results}

The following observations supplement the principal comparisons in
Chapter~\ref{sec:results}. Full scores, dispersion and paired effects
are retained in the tables above.


\subsubsection*{GeoFWI3D: consistency, near ties and negative outcomes}

All four terms give the highest mean for nnU-Net, SwinUNETR and the
compact U-Net, with improvements in all nine paired cells for each.
UNETR reaches its highest mean with $B+E$, improving eight of nine
cells and all three fold means. The compact U-Net has the highest
absolute overlap, while nnU-Net has the largest improvement over its
own baseline. For SwinUNETR, the all-term mean exceeds $B+S+D$ by
only 0.02 points. The compact U-Net's all-term and $B+S+D$ scores
both round to 80.73\%; their unrounded difference is below 0.01
points. These rankings do not establish meaningful superiority
between the near-tied combinations.

Boundary bands alone improve all four model means, by 1.35, 1.59,
0.75 and 0.72 points, respectively. However, only 30 of the 36 paired
cells improve, compared with all 72 medical cells. Depth alone also
has a positive mean effect for every geological model. These patterns
extend the evidence for annotation-derived supervision beyond MRI,
while showing that the consistency and size of the gains depend on
the task and recipe.

The full inventory retains negative findings. SDF alone changes
nnU-Net by $-0.07$ and UNETR by $-0.33$ points. nnU-Net's
$B+E+D$ mean is 1.36 points below baseline despite improving five
individual cells. UNETR's $B+E+S$ and $E+S+D$ also have negative
mean effects. Five constrained UNETR cells finish below 50\% Dice;
they remain in the summaries. No interpretation depends on discarding
those outcomes or replacing them with validation-best checkpoints.

\subsubsection*{MSD: individual constraints}

Boundary bands alone improve all four model means and all 36 matched
comparisons. Mean gains are 7.55, 3.85, 3.89 and 4.83 points for
nnU-Net, SwinUNETR, UNETR and the compact U-Net, respectively.


The inventory also contains negative findings. SDF alone reduces
nnU-Net by 2.04 points and is lower in all nine cells; its mean effect
on the compact U-Net is $-0.17$ points. Depth alone is positive for
all four model means, especially SwinUNETR (+8.24 points), but it is
not uniformly beneficial as an addition to another objective. Adding
edges to bands lowers the nnU-Net mean by 1.09 points while improving
the other three models. These are comparisons of calibrated configurations;
combining terms can change their coefficients as well as their domains.

\subsubsection*{ADNI: individual constraints}

Boundary bands alone improve every model mean and all 36 matched
comparisons. Gains are 0.73, 3.34, 2.99 and 3.81 points for nnU-Net,
SwinUNETR, UNETR and the compact U-Net, respectively.


SDF alone has essentially zero mean effect on nnU-Net (+0.01 points)
and negative mean effects on the other three models. Depth alone and
bands alone are positive for all model means. The preferred objective
is shared across datasets for nnU-Net and the compact U-Net, but differs
for the two transformer models. Neither absolute Dice nor gain should
be treated as a direct measure of relative dataset difficulty: MSD
requires two foreground labels, while ADNI has a single foreground,
different subjects and different crop geometry.

\subsubsection*{SAM3: variation across objectives and support partitions}

The adapted baseline reaches 69.31\% validation Dice and 74.97\%
support Dice. Adding bands, edges and depth supervision gives the
highest validation mean, 72.76\%, with a paired gain of
$3.45\pm2.43$ percentage points across the three fold means.
The gain is positive in eight of nine runs. Its fold-mean changes are
+6.19, +2.57 and +1.58 points for folds 0, 1 and 3, respectively
(Table~\ref{tab:sam3-geo-folds}). Thus the mean improvement occurs
in every support partition, although its magnitude varies and the
fold-0, seed-2 run loses 2.74 points. The selected configuration's
support Dice is 77.41\%; selection by the highest validation mean
remains exploratory.

Ten of the thirteen constrained variants improve mean validation Dice.
Depth supervision alone gives a 2.98-point gain, positive in eight
runs, while bands alone give 1.48 points, positive in six. The
$B+D$ and $B+S+D$ combinations also exceed the baseline by more
than 2.8 points. These results associate several of the larger gains
with depth supervision under the present calibration and schedule;
they do not isolate the contribution of depth weighting from that
of the added classification loss.

Combining more terms does not guarantee a larger benefit. Adding
signed distance to $B+E+D$ lowers mean Dice from 72.76\% to
70.95\%. Signed distance alone and $E+D$ give small negative
mean changes, while $B+E+S$ loses 1.85 points and improves only
three runs. The small positive means of $B+S$ and $B+E$ also
occur in only three runs each, illustrating why an average gain
should be read alongside its fold decomposition and positive-pair count.

\subsubsection*{MSD MedSAM3: matched execution batches}

The MSD MedSAM3 campaign evaluates fourteen objectives across three
folds and three seeds, with five labelled support crops and 200 updates
per run. Table~\ref{tab:medsam-msd-noaug} presents the complete set of
constrained variants together with their matched baseline scores.
The controls for $B$ and $B+E$ reach $74.93\pm2.18$\% validation Dice
and 79.18\% training-set Dice; those for the other variants reach
$76.34\pm1.64$\% and 81.00\%, respectively. Gains retain these pairings.


Baseline $+B$ reaches 78.76\% mean validation A/P Dice, a 3.83-point improvement
with eight positive pairs out of nine. Baseline $+B+E$ reaches 80.12\%,
improving its matched baseline by 5.19 points in all nine pairs.
Adding edges to bands increases the mean by 1.36 points within this
matched comparison.

The highest mean is 81.73\% with Baseline $+B+E+S+D$, a 5.38-point
gain over its 76.34\% control. The three fold-mean gains are 4.74,
5.71 and 5.71 points, and all nine paired differences are positive.
Baseline $+S$ gives a smaller +0.62-point mean change with six positive
pairs; Baseline $+D$ gives +2.56 points with eight. The available
combinations demonstrate benefits beyond those single-family variants.

The campaign was executed in batches with distinct calibration records
and execution environments. Comparisons across batches therefore
describe the achieved scores and paired gains; a difference between two
such variants cannot isolate the effect of adding a constraint.
All fourteen objectives have completed results.

\subsubsection*{ADNI MedSAM3: closely ranked configurations}

On ADNI, $E+S+D$ has the highest mean, 80.85\%, compared with 77.32\%
for baseline. Its +3.53-point gain is positive in all nine cells;
fold-mean gains are 3.36, 3.77 and 3.46 points. The all-family
configuration reaches 80.74\%. Their 0.11-point separation is a
descriptive ranking, not evidence of a reliable superiority of one
combination. Bands alone improve the mean by 1.75 points with eight
positive pairs; $B+E$ improves it by 1.91 points with seven. The
shared support crops limit inference about sensitivity to which two
annotated examples are chosen.

\FloatBarrier

\subsection{Complete Error Counts}\label{sec:complete-error-tables}

\label{sec:error-details}


Higher Dice can result from recovering more of an object, reducing
foreground outside it, or both. A false negative (FN) is foreground
predicted as background; a false positive (FP) is background predicted as
foreground. MSD additionally distinguishes anterior/posterior swaps,
which preserve whole-hippocampus membership but assign the wrong internal
class. These categories count different mistakes, so improved overlap
does not imply that every error type decreases.

For the conventional geological models, the highest-mean nnU-Net objective
reduces both FN (34.4\%) and FP (11.4\%). SwinUNETR and the compact U-Net
mainly reduce FP, by 12.4\% and 26.9\%, while their FN counts rise
slightly. UNETR shows the reverse tradeoff: FN falls by 4.3\% and FP
rises by 2.9\%. Thus, even within one dataset, spatial supervision changes
foreground retention and extension differently across architectures.

The conventional medical results show a related pattern. On MSD, the
selected nnU-Net objective reduces missed foreground with a slight
increase in extension; the other models reduce extension while increasing
omissions. All four selected configurations reduce anterior/posterior
swaps. On ADNI, nnU-Net reduces both binary error types, while the other
models exchange fewer false positives for more false negatives.

Pretrained adaptation favours greater foreground recovery. SAM3's
highest-mean objective reduces FN from 34.16 to 27.16 thousand voxels
per volume, while FP rises from 3.70 to 6.99 thousand. Total incorrectly
labelled voxels nevertheless fall by 9.8\%. For MedSAM3 on MSD, the
selected objective reduces FN by 36.0\% and internal swaps by 20.7\%,
while FP rises by 40.7\%. On ADNI, FN falls by 24.0\% and FP rises
by 26.2\%. These changes explain why an overlap gain should be considered
together with the error balance relevant to the application.

All counts cover the full evaluated image, including locations beyond
the constraint bands. Tables~\ref{tab:geo-all-errors}--\ref{tab:medsam-adni-all-errors}
below report every objective, including
those with negative Dice changes. Conventional-model counts are pooled
within each run; adaptation counts are means per image. Their absolute
magnitudes therefore should not be compared across studies. The error
counts establish how much misclassification changes. The separate
qualitative reproductions in Section~\ref{sec:qualitative-predictions}
provide matched prediction examples for the pretrained models; they do
not replace the original aggregate error counts below. Surface-distance
measurements remain outside the present evaluation.

\FloatBarrier

Counts include all evaluated voxels. Conventional-model counts pool the
validation images within each run; adaptation counts are means per image.
Each condition then averages seeds within folds and weights folds equally.
Dispersion is the sample standard deviation across the three fold means.
Both the mean and its dispersion use the scale in each column heading.
Rows ascend by mean total error within a model, using unrounded values.

\begingroup\footnotesize\setstretch{1.0}
\setlength{\tabcolsep}{4pt}
\renewcommand{\arraystretch}{1.12}
\begin{longtable}{@{}lrr@{}}
\caption{GeoFWI3D conventional models: all fourteen objectives for each architecture. Entries are mean $\pm$ sample SD across three seed-averaged fold means, using all nine runs per row. Counts pool the full validation set within each run. Within each model, rows ascend by mean total error (FN $+$ FP).}
\label{tab:geo-all-errors}\\
\hline
Variant & \shortstack{FN / run\\($\times 10^{6}$)} & \shortstack{FP / run\\($\times 10^{6}$)} \\ \hline
\endfirsthead
\hline
Variant & \shortstack{FN / run\\($\times 10^{6}$)} & \shortstack{FP / run\\($\times 10^{6}$)} \\ \hline
\endhead
\hline
\multicolumn{3}{r}{Continued on next page}\\
\endfoot
\hline\endlastfoot
\multicolumn{3}{@{}l}{\textbf{nnU-Net}}\\*
Baseline $+B+E+S+D$ & 1.39 $\pm$ 0.57 & 1.21 $\pm$ 0.52 \\
Baseline $+B+S+D$ & 1.50 $\pm$ 0.53 & 1.22 $\pm$ 0.52 \\
Baseline $+E+S+D$ & 1.54 $\pm$ 0.58 & 1.25 $\pm$ 0.66 \\
Baseline $+S+D$ & 1.74 $\pm$ 0.65 & 1.15 $\pm$ 0.54 \\
Baseline $+B+D$ & 1.52 $\pm$ 0.56 & 1.47 $\pm$ 0.67 \\
Baseline $+E+D$ & 1.66 $\pm$ 0.58 & 1.35 $\pm$ 0.61 \\
Baseline $+D$ & 2.05 $\pm$ 0.76 & 1.18 $\pm$ 0.49 \\
Baseline $+B$ & 1.97 $\pm$ 0.78 & 1.31 $\pm$ 0.53 \\
Baseline $+B+S$ & 2.21 $\pm$ 0.79 & 1.11 $\pm$ 0.46 \\
Baseline $+B+E$ & 1.85 $\pm$ 0.78 & 1.53 $\pm$ 0.66 \\
Baseline $+B+E+S$ & 2.31 $\pm$ 0.95 & 1.13 $\pm$ 0.46 \\
Baseline $+S$ & 2.16 $\pm$ 0.70 & 1.33 $\pm$ 0.55 \\
Baseline & 2.12 $\pm$ 0.70 & 1.36 $\pm$ 0.57 \\
Baseline $+B+E+D$ & 1.05 $\pm$ 0.26 & 2.73 $\pm$ 0.67 \\
\hline
\multicolumn{3}{@{}l}{\textbf{SwinUNETR}}\\*
Baseline $+B+E+S+D$ & 2.14 $\pm$ 0.77 & 4.40 $\pm$ 1.08 \\
Baseline $+B+S+D$ & 2.13 $\pm$ 0.77 & 4.42 $\pm$ 1.10 \\
Baseline $+B+D$ & 2.15 $\pm$ 0.79 & 4.45 $\pm$ 1.10 \\
Baseline $+B+E+D$ & 2.15 $\pm$ 0.75 & 4.46 $\pm$ 1.12 \\
Baseline $+B+E$ & 2.11 $\pm$ 0.73 & 4.62 $\pm$ 1.06 \\
Baseline $+B$ & 2.08 $\pm$ 0.71 & 4.66 $\pm$ 1.04 \\
Baseline $+B+S$ & 2.11 $\pm$ 0.73 & 4.64 $\pm$ 1.08 \\
Baseline $+B+E+S$ & 2.10 $\pm$ 0.74 & 4.67 $\pm$ 1.10 \\
Baseline $+E+S+D$ & 2.14 $\pm$ 0.76 & 4.64 $\pm$ 0.98 \\
Baseline $+E+D$ & 2.13 $\pm$ 0.74 & 4.67 $\pm$ 0.98 \\
Baseline $+D$ & 2.15 $\pm$ 0.75 & 4.70 $\pm$ 0.95 \\
Baseline $+S+D$ & 2.16 $\pm$ 0.80 & 4.69 $\pm$ 1.05 \\
Baseline $+S$ & 2.10 $\pm$ 0.79 & 4.99 $\pm$ 0.82 \\
Baseline & 2.13 $\pm$ 0.82 & 5.02 $\pm$ 0.90 \\
\hline
\multicolumn{3}{@{}l}{\textbf{UNETR}}\\*
Baseline $+B+E+S+D$ & 3.67 $\pm$ 1.06 & 1.74 $\pm$ 0.82 \\
Baseline $+B+S+D$ & 3.66 $\pm$ 1.04 & 1.81 $\pm$ 0.84 \\
Baseline $+S+D$ & 3.71 $\pm$ 1.02 & 1.77 $\pm$ 0.78 \\
Baseline $+E+D$ & 3.51 $\pm$ 1.02 & 1.99 $\pm$ 0.85 \\
Baseline $+B+S$ & 3.70 $\pm$ 1.03 & 1.80 $\pm$ 0.76 \\
Baseline $+E+S+D$ & 3.76 $\pm$ 1.04 & 1.75 $\pm$ 0.83 \\
Baseline $+B+E+D$ & 3.42 $\pm$ 0.98 & 2.09 $\pm$ 0.89 \\
Baseline $+B+E+S$ & 3.77 $\pm$ 1.05 & 1.74 $\pm$ 0.79 \\
Baseline $+D$ & 3.49 $\pm$ 1.02 & 2.03 $\pm$ 0.87 \\
Baseline $+B+E$ & 3.42 $\pm$ 1.03 & 2.11 $\pm$ 0.90 \\
Baseline $+B+D$ & 3.45 $\pm$ 0.99 & 2.09 $\pm$ 0.91 \\
Baseline $+B$ & 3.46 $\pm$ 0.98 & 2.12 $\pm$ 0.89 \\
Baseline $+S$ & 3.73 $\pm$ 1.03 & 1.86 $\pm$ 0.77 \\
Baseline & 3.58 $\pm$ 0.97 & 2.05 $\pm$ 0.78 \\
\hline
\multicolumn{3}{@{}l}{\textbf{3D U-Net}}\\*
Baseline $+B+E+S+D$ & 0.72 $\pm$ 0.31 & 1.44 $\pm$ 0.45 \\
Baseline $+B+S+D$ & 0.71 $\pm$ 0.31 & 1.46 $\pm$ 0.47 \\
Baseline $+B+E+D$ & 0.73 $\pm$ 0.32 & 1.46 $\pm$ 0.45 \\
Baseline $+B+D$ & 0.70 $\pm$ 0.31 & 1.50 $\pm$ 0.49 \\
Baseline $+E+S+D$ & 0.67 $\pm$ 0.30 & 1.61 $\pm$ 0.52 \\
Baseline $+S+D$ & 0.69 $\pm$ 0.30 & 1.67 $\pm$ 0.56 \\
Baseline $+E+D$ & 0.67 $\pm$ 0.29 & 1.72 $\pm$ 0.59 \\
Baseline $+D$ & 0.66 $\pm$ 0.31 & 1.76 $\pm$ 0.65 \\
Baseline $+B+E+S$ & 0.70 $\pm$ 0.29 & 1.79 $\pm$ 0.46 \\
Baseline $+B$ & 0.68 $\pm$ 0.29 & 1.85 $\pm$ 0.55 \\
Baseline $+B+E$ & 0.68 $\pm$ 0.27 & 1.85 $\pm$ 0.40 \\
Baseline $+B+S$ & 0.68 $\pm$ 0.27 & 1.93 $\pm$ 0.33 \\
Baseline $+S$ & 0.71 $\pm$ 0.31 & 1.93 $\pm$ 0.60 \\
Baseline & 0.69 $\pm$ 0.29 & 1.97 $\pm$ 0.58 \\
\hline
\end{longtable}
\endgroup

\begingroup\footnotesize\setstretch{1.0}
\setlength{\tabcolsep}{4pt}
\renewcommand{\arraystretch}{1.12}
\begin{longtable}{@{}lrrr@{}}
\caption{MSD conventional models: all fourteen objectives for each architecture. Entries are mean $\pm$ sample SD across three seed-averaged fold means, using all nine runs per row. Counts pool the full validation set within each run. Within each model, rows ascend by mean total error (FN $+$ FP $+$ A/P swaps).}
\label{tab:msd-all-errors}\label{tab:capped-errors}\\
\hline
Variant & \shortstack{FN / run\\($\times 10^{3}$)} & \shortstack{FP / run\\($\times 10^{3}$)} & \shortstack{A/P swaps / run\\($\times 10^{3}$)} \\ \hline
\endfirsthead
\hline
Variant & \shortstack{FN / run\\($\times 10^{3}$)} & \shortstack{FP / run\\($\times 10^{3}$)} & \shortstack{A/P swaps / run\\($\times 10^{3}$)} \\ \hline
\endhead
\hline
\multicolumn{4}{r}{Continued on next page}\\
\endfoot
\hline\endlastfoot
\multicolumn{4}{@{}l}{\textbf{nnU-Net}}\\*
Baseline $+B+E+S+D$ & 34.63 $\pm$ 4.36 & 29.14 $\pm$ 1.68 & 12.77 $\pm$ 1.71 \\
Baseline $+B+S+D$ & 34.09 $\pm$ 3.55 & 30.26 $\pm$ 1.51 & 12.60 $\pm$ 1.72 \\
Baseline $+B+E+D$ & 35.11 $\pm$ 3.94 & 29.71 $\pm$ 1.44 & 13.62 $\pm$ 1.82 \\
Baseline $+B+D$ & 34.40 $\pm$ 3.57 & 31.15 $\pm$ 1.45 & 13.29 $\pm$ 1.78 \\
Baseline $+B+S$ & 37.49 $\pm$ 4.53 & 30.74 $\pm$ 1.66 & 13.60 $\pm$ 1.53 \\
Baseline $+B$ & 36.75 $\pm$ 4.52 & 31.58 $\pm$ 2.08 & 14.06 $\pm$ 1.60 \\
Baseline $+B+E+S$ & 39.62 $\pm$ 5.18 & 28.92 $\pm$ 2.24 & 14.16 $\pm$ 1.55 \\
Baseline $+B+E$ & 38.75 $\pm$ 5.19 & 29.67 $\pm$ 2.30 & 14.82 $\pm$ 1.72 \\
Baseline $+S+D$ & 48.03 $\pm$ 9.09 & 27.32 $\pm$ 3.95 & 12.86 $\pm$ 1.49 \\
Baseline $+E+S+D$ & 49.08 $\pm$ 10.10 & 26.00 $\pm$ 4.31 & 13.19 $\pm$ 1.56 \\
Baseline $+D$ & 49.87 $\pm$ 10.13 & 27.43 $\pm$ 4.69 & 13.49 $\pm$ 1.57 \\
Baseline $+E+D$ & 51.49 $\pm$ 11.02 & 26.05 $\pm$ 4.88 & 14.03 $\pm$ 1.65 \\
Baseline & 51.25 $\pm$ 12.54 & 28.97 $\pm$ 6.06 & 14.12 $\pm$ 1.68 \\
Baseline $+S$ & 56.56 $\pm$ 13.76 & 26.18 $\pm$ 5.60 & 13.42 $\pm$ 1.08 \\
\hline
\multicolumn{4}{@{}l}{\textbf{SwinUNETR}}\\*
Baseline $+B+E+S+D$ & 5.36 $\pm$ 1.28 & 155.79 $\pm$ 6.00 & 5.96 $\pm$ 0.77 \\
Baseline $+B+S+D$ & 4.96 $\pm$ 1.18 & 158.26 $\pm$ 5.90 & 5.88 $\pm$ 0.70 \\
Baseline $+B+E+D$ & 5.85 $\pm$ 1.38 & 160.75 $\pm$ 5.84 & 5.95 $\pm$ 0.75 \\
Baseline $+B+D$ & 5.10 $\pm$ 1.40 & 165.24 $\pm$ 7.01 & 5.85 $\pm$ 0.64 \\
Baseline $+E+S+D$ & 2.79 $\pm$ 0.84 & 178.30 $\pm$ 9.94 & 5.98 $\pm$ 0.50 \\
Baseline $+E+D$ & 2.92 $\pm$ 0.87 & 180.22 $\pm$ 10.23 & 5.97 $\pm$ 0.52 \\
Baseline $+S+D$ & 2.54 $\pm$ 0.70 & 183.92 $\pm$ 10.78 & 6.06 $\pm$ 0.57 \\
Baseline $+D$ & 2.61 $\pm$ 0.72 & 186.21 $\pm$ 9.53 & 6.02 $\pm$ 0.49 \\
Baseline $+B+S$ & 4.27 $\pm$ 1.06 & 293.26 $\pm$ 63.55 & 5.98 $\pm$ 0.56 \\
Baseline $+B+E+S$ & 4.61 $\pm$ 1.07 & 295.88 $\pm$ 61.75 & 6.00 $\pm$ 0.45 \\
Baseline $+B$ & 4.42 $\pm$ 0.93 & 324.20 $\pm$ 68.15 & 5.95 $\pm$ 0.42 \\
Baseline $+B+E$ & 5.12 $\pm$ 1.23 & 326.80 $\pm$ 63.71 & 6.03 $\pm$ 0.44 \\
Baseline $+S$ & 1.45 $\pm$ 0.42 & 345.53 $\pm$ 49.45 & 6.21 $\pm$ 0.27 \\
Baseline & 1.49 $\pm$ 0.41 & 388.34 $\pm$ 55.21 & 6.36 $\pm$ 0.29 \\
\hline
\multicolumn{4}{@{}l}{\textbf{UNETR}}\\*
Baseline $+B+E+S$ & 12.27 $\pm$ 2.96 & 206.13 $\pm$ 11.59 & 6.18 $\pm$ 0.59 \\
Baseline $+B+S$ & 11.79 $\pm$ 3.04 & 208.42 $\pm$ 17.42 & 6.21 $\pm$ 0.56 \\
Baseline $+B+E+S+D$ & 13.89 $\pm$ 3.43 & 213.83 $\pm$ 19.18 & 6.01 $\pm$ 0.47 \\
Baseline $+B+S+D$ & 13.12 $\pm$ 3.00 & 216.21 $\pm$ 19.88 & 6.07 $\pm$ 0.63 \\
Baseline $+B+E$ & 12.62 $\pm$ 3.15 & 225.02 $\pm$ 19.89 & 6.26 $\pm$ 0.63 \\
Baseline $+B$ & 12.06 $\pm$ 2.89 & 227.37 $\pm$ 21.72 & 6.23 $\pm$ 0.57 \\
Baseline $+B+E+D$ & 14.14 $\pm$ 3.27 & 230.27 $\pm$ 24.66 & 6.09 $\pm$ 0.62 \\
Baseline $+B+D$ & 13.65 $\pm$ 3.41 & 232.29 $\pm$ 27.12 & 6.05 $\pm$ 0.66 \\
Baseline $+E+S+D$ & 9.03 $\pm$ 2.80 & 237.84 $\pm$ 20.07 & 6.40 $\pm$ 0.70 \\
Baseline $+S+D$ & 8.68 $\pm$ 2.78 & 241.21 $\pm$ 18.68 & 6.39 $\pm$ 0.64 \\
Baseline $+E+D$ & 9.43 $\pm$ 2.92 & 247.54 $\pm$ 24.11 & 6.33 $\pm$ 0.63 \\
Baseline $+D$ & 9.02 $\pm$ 2.82 & 252.46 $\pm$ 24.46 & 6.39 $\pm$ 0.67 \\
Baseline $+S$ & 6.32 $\pm$ 1.90 & 255.86 $\pm$ 10.53 & 6.60 $\pm$ 0.58 \\
Baseline & 6.74 $\pm$ 2.27 & 270.61 $\pm$ 15.13 & 6.73 $\pm$ 0.76 \\
\hline
\multicolumn{4}{@{}l}{\textbf{3D U-Net}}\\*
Baseline $+B+E+D$ & 10.26 $\pm$ 1.54 & 57.47 $\pm$ 4.99 & 4.86 $\pm$ 0.30 \\
Baseline $+B+E+S+D$ & 10.23 $\pm$ 1.63 & 58.02 $\pm$ 6.63 & 4.89 $\pm$ 0.41 \\
Baseline $+B+E$ & 9.78 $\pm$ 1.34 & 60.30 $\pm$ 5.41 & 4.92 $\pm$ 0.49 \\
Baseline $+B+E+S$ & 9.46 $\pm$ 1.17 & 61.43 $\pm$ 5.00 & 5.05 $\pm$ 0.29 \\
Baseline $+B+D$ & 9.10 $\pm$ 1.49 & 64.49 $\pm$ 6.53 & 5.05 $\pm$ 0.25 \\
Baseline $+B+S+D$ & 8.79 $\pm$ 1.57 & 65.31 $\pm$ 7.46 & 5.06 $\pm$ 0.45 \\
Baseline $+B$ & 8.67 $\pm$ 1.35 & 66.06 $\pm$ 7.32 & 5.07 $\pm$ 0.33 \\
Baseline $+B+S$ & 8.76 $\pm$ 1.38 & 66.17 $\pm$ 7.05 & 4.95 $\pm$ 0.28 \\
Baseline $+E+S+D$ & 7.49 $\pm$ 1.38 & 75.95 $\pm$ 8.27 & 5.02 $\pm$ 0.57 \\
Baseline $+E+D$ & 6.80 $\pm$ 0.88 & 77.32 $\pm$ 6.69 & 5.14 $\pm$ 0.52 \\
Baseline $+S+D$ & 6.51 $\pm$ 0.82 & 82.57 $\pm$ 6.71 & 5.09 $\pm$ 0.32 \\
Baseline $+D$ & 6.47 $\pm$ 0.96 & 82.89 $\pm$ 8.81 & 5.12 $\pm$ 0.39 \\
Baseline & 4.96 $\pm$ 0.68 & 99.43 $\pm$ 8.64 & 4.99 $\pm$ 0.15 \\
Baseline $+S$ & 5.21 $\pm$ 1.18 & 99.03 $\pm$ 11.06 & 5.32 $\pm$ 0.61 \\
\hline
\end{longtable}
\endgroup

\begingroup\footnotesize\setstretch{1.0}
\setlength{\tabcolsep}{4pt}
\renewcommand{\arraystretch}{1.12}
\begin{longtable}{@{}lrr@{}}
\caption{ADNI conventional models: all fourteen objectives for each architecture. Entries are mean $\pm$ sample SD across three seed-averaged fold means, using all nine runs per row. Counts pool the full validation set within each run. Within each model, rows ascend by mean total error (FN $+$ FP).}
\label{tab:adni-all-errors}\\
\hline
Variant & \shortstack{FN / run\\($\times 10^{3}$)} & \shortstack{FP / run\\($\times 10^{3}$)} \\ \hline
\endfirsthead
\hline
Variant & \shortstack{FN / run\\($\times 10^{3}$)} & \shortstack{FP / run\\($\times 10^{3}$)} \\ \hline
\endhead
\hline
\multicolumn{3}{r}{Continued on next page}\\
\endfoot
\hline\endlastfoot
\multicolumn{3}{@{}l}{\textbf{nnU-Net}}\\*
Baseline $+B+S+D$ & 35.91 $\pm$ 9.08 & 32.26 $\pm$ 10.23 \\
Baseline $+B+D$ & 35.94 $\pm$ 9.22 & 32.48 $\pm$ 10.60 \\
Baseline $+B+E+S+D$ & 36.29 $\pm$ 9.63 & 32.26 $\pm$ 11.28 \\
Baseline $+B+E+D$ & 36.03 $\pm$ 9.74 & 32.76 $\pm$ 11.47 \\
Baseline $+B+S$ & 36.01 $\pm$ 9.61 & 33.27 $\pm$ 11.00 \\
Baseline $+B$ & 36.03 $\pm$ 9.40 & 33.25 $\pm$ 10.95 \\
Baseline $+E+S+D$ & 37.25 $\pm$ 9.95 & 32.07 $\pm$ 11.22 \\
Baseline $+E+D$ & 37.02 $\pm$ 9.84 & 32.63 $\pm$ 11.42 \\
Baseline $+B+E+S$ & 36.10 $\pm$ 10.19 & 33.66 $\pm$ 12.17 \\
Baseline $+S+D$ & 36.94 $\pm$ 9.74 & 32.93 $\pm$ 11.42 \\
Baseline $+B+E$ & 35.97 $\pm$ 10.14 & 33.91 $\pm$ 12.32 \\
Baseline $+D$ & 36.85 $\pm$ 9.86 & 33.54 $\pm$ 11.73 \\
Baseline $+S$ & 36.94 $\pm$ 10.12 & 34.58 $\pm$ 12.83 \\
Baseline & 36.91 $\pm$ 10.16 & 34.68 $\pm$ 12.84 \\
\hline
\multicolumn{3}{@{}l}{\textbf{SwinUNETR}}\\*
Baseline $+B+E+D$ & 16.89 $\pm$ 2.41 & 78.63 $\pm$ 9.19 \\
Baseline $+B+E+S+D$ & 16.89 $\pm$ 2.44 & 79.05 $\pm$ 9.99 \\
Baseline $+B+D$ & 16.21 $\pm$ 2.48 & 81.71 $\pm$ 9.57 \\
Baseline $+B+S+D$ & 16.29 $\pm$ 2.34 & 82.12 $\pm$ 10.53 \\
Baseline $+B+E$ & 15.86 $\pm$ 2.13 & 83.41 $\pm$ 9.75 \\
Baseline $+B+E+S$ & 15.99 $\pm$ 2.36 & 83.48 $\pm$ 10.43 \\
Baseline $+B$ & 15.35 $\pm$ 2.19 & 85.97 $\pm$ 10.17 \\
Baseline $+B+S$ & 15.60 $\pm$ 2.31 & 85.89 $\pm$ 11.03 \\
Baseline $+E+D$ & 14.58 $\pm$ 2.42 & 95.41 $\pm$ 12.06 \\
Baseline $+E+S+D$ & 14.57 $\pm$ 2.58 & 95.54 $\pm$ 12.12 \\
Baseline $+D$ & 14.27 $\pm$ 2.34 & 98.21 $\pm$ 11.31 \\
Baseline $+S+D$ & 14.11 $\pm$ 2.40 & 98.72 $\pm$ 12.25 \\
Baseline & 13.17 $\pm$ 2.94 & 107.73 $\pm$ 13.17 \\
Baseline $+S$ & 13.26 $\pm$ 2.83 & 107.95 $\pm$ 11.99 \\
\hline
\multicolumn{3}{@{}l}{\textbf{UNETR}}\\*
Baseline $+B+E+S+D$ & 25.22 $\pm$ 5.03 & 140.23 $\pm$ 9.58 \\
Baseline $+B+E+D$ & 25.61 $\pm$ 5.30 & 140.71 $\pm$ 11.95 \\
Baseline $+B+E+S$ & 23.79 $\pm$ 4.88 & 145.42 $\pm$ 15.00 \\
Baseline $+B+D$ & 24.01 $\pm$ 4.86 & 147.25 $\pm$ 13.04 \\
Baseline $+B+S+D$ & 23.48 $\pm$ 4.92 & 149.29 $\pm$ 12.88 \\
Baseline $+B+E$ & 23.96 $\pm$ 4.95 & 149.13 $\pm$ 17.16 \\
Baseline $+B+S$ & 22.48 $\pm$ 4.20 & 153.58 $\pm$ 14.02 \\
Baseline $+B$ & 22.59 $\pm$ 4.94 & 156.09 $\pm$ 15.65 \\
Baseline $+E+S+D$ & 20.99 $\pm$ 4.71 & 166.67 $\pm$ 17.18 \\
Baseline $+E+D$ & 21.51 $\pm$ 4.73 & 168.22 $\pm$ 18.36 \\
Baseline $+S+D$ & 19.70 $\pm$ 4.40 & 174.31 $\pm$ 13.34 \\
Baseline $+D$ & 20.16 $\pm$ 4.62 & 174.23 $\pm$ 14.84 \\
Baseline & 17.26 $\pm$ 4.09 & 192.93 $\pm$ 18.30 \\
Baseline $+S$ & 16.96 $\pm$ 3.78 & 196.19 $\pm$ 17.67 \\
\hline
\multicolumn{3}{@{}l}{\textbf{3D U-Net}}\\*
Baseline $+B+E+D$ & 17.35 $\pm$ 3.67 & 67.90 $\pm$ 10.34 \\
Baseline $+B+E+S+D$ & 16.98 $\pm$ 3.84 & 69.22 $\pm$ 11.68 \\
Baseline $+B+E+S$ & 16.15 $\pm$ 3.58 & 73.55 $\pm$ 14.26 \\
Baseline $+B+S+D$ & 16.17 $\pm$ 4.36 & 73.63 $\pm$ 12.94 \\
Baseline $+B+S$ & 15.93 $\pm$ 3.64 & 74.21 $\pm$ 9.93 \\
Baseline $+B+E$ & 16.66 $\pm$ 3.76 & 73.76 $\pm$ 15.96 \\
Baseline $+B+D$ & 15.61 $\pm$ 4.27 & 75.50 $\pm$ 11.53 \\
Baseline $+B$ & 15.57 $\pm$ 3.66 & 78.91 $\pm$ 8.31 \\
Baseline $+E+S+D$ & 15.75 $\pm$ 3.49 & 83.92 $\pm$ 12.91 \\
Baseline $+E+D$ & 15.13 $\pm$ 4.70 & 87.13 $\pm$ 11.78 \\
Baseline $+D$ & 14.41 $\pm$ 3.28 & 89.93 $\pm$ 8.80 \\
Baseline $+S+D$ & 14.17 $\pm$ 3.30 & 90.28 $\pm$ 9.46 \\
Baseline & 12.91 $\pm$ 3.24 & 103.28 $\pm$ 14.28 \\
Baseline $+S$ & 12.64 $\pm$ 3.70 & 105.25 $\pm$ 14.81 \\
\hline
\end{longtable}
\endgroup

\begingroup\footnotesize\setstretch{1.0}
\setlength{\tabcolsep}{4pt}
\renewcommand{\arraystretch}{1.12}
\begin{longtable}{@{}lrrr@{}}
\caption{SAM3 on GeoFWI3D: full-volume FN and FP for all fourteen variants. Each run averages errors over 96 held-out volumes; rows report mean $\pm$ sample SD across three seed-averaged fold means. All rows contain nine runs. Rows ascend by mean total error (FN $+$ FP).}
\label{tab:sam3-geo-errors}\\
\hline
Variant & \shortstack{FN / volume\\($\times 10^{3}$)} & \shortstack{FP / volume\\($\times 10^{3}$)} & Runs \\ \hline
\endfirsthead
\hline
Variant & \shortstack{FN / volume\\($\times 10^{3}$)} & \shortstack{FP / volume\\($\times 10^{3}$)} & Runs \\ \hline
\endhead
\hline
\multicolumn{4}{r}{Continued on next page}\\
\endfoot
\hline\endlastfoot
Baseline $+B+E+D$ & 27.16 $\pm$ 5.11 & 6.99 $\pm$ 1.69 & 9/9 \\
Baseline $+B+D$ & 28.64 $\pm$ 3.16 & 5.81 $\pm$ 1.70 & 9/9 \\
Baseline $+B+S+D$ & 29.29 $\pm$ 3.73 & 5.59 $\pm$ 1.92 & 9/9 \\
Baseline $+D$ & 30.19 $\pm$ 4.29 & 4.80 $\pm$ 1.18 & 9/9 \\
Baseline $+B+E+S+D$ & 30.58 $\pm$ 5.39 & 5.41 $\pm$ 2.60 & 9/9 \\
Baseline $+B$ & 29.74 $\pm$ 5.03 & 6.38 $\pm$ 3.08 & 9/9 \\
Baseline $+E+S+D$ & 32.84 $\pm$ 5.64 & 4.40 $\pm$ 1.16 & 9/9 \\
Baseline $+S+D$ & 32.19 $\pm$ 3.24 & 5.15 $\pm$ 1.58 & 9/9 \\
Baseline $+B+E$ & 32.46 $\pm$ 7.51 & 5.12 $\pm$ 2.59 & 9/9 \\
Baseline $+S$ & 33.82 $\pm$ 6.83 & 3.97 $\pm$ 1.40 & 9/9 \\
Baseline & 34.16 $\pm$ 6.41 & 3.70 $\pm$ 0.49 & 9/9 \\
Baseline $+E+D$ & 32.76 $\pm$ 6.28 & 5.16 $\pm$ 2.13 & 9/9 \\
Baseline $+B+S$ & 33.94 $\pm$ 3.05 & 4.12 $\pm$ 1.25 & 9/9 \\
Baseline $+B+E+S$ & 36.05 $\pm$ 4.36 & 3.87 $\pm$ 1.48 & 9/9 \\
\end{longtable}
\endgroup

\begingroup\footnotesize\setstretch{1.0}
\setlength{\tabcolsep}{4pt}
\renewcommand{\arraystretch}{1.12}
\begin{longtable}{@{}lrrr@{}}
\caption{MSD MedSAM3: all variants and their matched controls. Counts are mean errors per validation crop within a run, then averaged over seeds and folds; $\pm$ gives sample SD across the three fold means. Each row includes nine runs. The two baseline labels identify their matched constrained variants. Within each model, rows ascend by mean total error (FN $+$ FP $+$ A/P swaps).}
\label{tab:medsam-msd-all-errors}\label{tab:medsam-noaug-errors}\label{tab:medsam-noaug-error-ledger}\\
\hline
Variant & \shortstack{FN / crop\\($\times 10^{2}$)} & \shortstack{FP / crop\\($\times 10^{2}$)} & \shortstack{A/P swaps / crop\\($\times 10^{2}$)} \\ \hline
\endfirsthead
\hline
Variant & \shortstack{FN / crop\\($\times 10^{2}$)} & \shortstack{FP / crop\\($\times 10^{2}$)} & \shortstack{A/P swaps / crop\\($\times 10^{2}$)} \\ \hline
\endhead
\hline
\multicolumn{4}{r}{Continued on next page}\\
\endfoot
\hline\endlastfoot
Baseline $+B+E+S+D$ & 5.56 $\pm$ 0.59 & 4.37 $\pm$ 0.22 & 0.98 $\pm$ 0.17 \\
Baseline $+E+S+D$ & 6.30 $\pm$ 0.66 & 3.73 $\pm$ 0.08 & 0.93 $\pm$ 0.19 \\
Baseline $+B+S+D$ & 6.27 $\pm$ 0.71 & 4.01 $\pm$ 0.22 & 0.91 $\pm$ 0.13 \\
Baseline $+E+D$ & 6.33 $\pm$ 0.70 & 4.02 $\pm$ 0.22 & 0.98 $\pm$ 0.24 \\
Baseline $+B+S$ & 6.77 $\pm$ 0.45 & 3.68 $\pm$ 0.28 & 0.90 $\pm$ 0.16 \\
Baseline $+B+E+D$ & 5.59 $\pm$ 0.57 & 4.75 $\pm$ 0.38 & 1.07 $\pm$ 0.17 \\
Baseline $+B+E+S$ & 5.73 $\pm$ 0.53 & 4.70 $\pm$ 0.39 & 0.99 $\pm$ 0.22 \\
Baseline $+S+D$ & 7.79 $\pm$ 0.48 & 2.95 $\pm$ 0.25 & 0.81 $\pm$ 0.13 \\
Baseline $+B+D$ & 6.32 $\pm$ 0.67 & 4.25 $\pm$ 0.08 & 0.99 $\pm$ 0.12 \\
Baseline $+D$ & 7.87 $\pm$ 0.59 & 3.14 $\pm$ 0.12 & 0.96 $\pm$ 0.32 \\
Baseline $+B+E$ & 5.55 $\pm$ 0.77 & 5.64 $\pm$ 0.38 & 1.04 $\pm$ 0.13 \\
Baseline $+B$ & 6.27 $\pm$ 0.66 & 5.16 $\pm$ 0.44 & 1.23 $\pm$ 0.08 \\
Baseline $+S$ & 9.54 $\pm$ 1.03 & 2.34 $\pm$ 0.24 & 0.80 $\pm$ 0.26 \\
Baseline (other variants) & 8.69 $\pm$ 0.83 & 3.11 $\pm$ 0.13 & 1.23 $\pm$ 0.38 \\
Baseline ($B$, $B+E$) & 9.00 $\pm$ 0.82 & 4.23 $\pm$ 1.19 & 1.11 $\pm$ 0.15 \\
\end{longtable}
\endgroup

\begingroup\footnotesize\setstretch{1.0}
\setlength{\tabcolsep}{4pt}
\renewcommand{\arraystretch}{1.12}
\begin{longtable}{@{}lrr@{}}
\caption{ADNI MedSAM3: all fourteen variants, each with nine runs. Counts are mean errors per validation crop within each run; entries give mean $\pm$ sample SD across three seed-averaged fold means. Within each model, rows ascend by mean total error (FN $+$ FP).}
\label{tab:medsam-adni-all-errors}\\
\hline
Variant & \shortstack{FN / crop\\($\times 10^{2}$)} & \shortstack{FP / crop\\($\times 10^{2}$)} \\ \hline
\endfirsthead
\hline
Variant & \shortstack{FN / crop\\($\times 10^{2}$)} & \shortstack{FP / crop\\($\times 10^{2}$)} \\ \hline
\endhead
\hline
\multicolumn{3}{r}{Continued on next page}\\
\endfoot
\hline\endlastfoot
Baseline $+E+S+D$ & 5.91 $\pm$ 0.37 & 3.59 $\pm$ 0.37 \\
Baseline $+B+S$ & 6.31 $\pm$ 0.40 & 3.26 $\pm$ 0.56 \\
Baseline $+B+E+S$ & 6.02 $\pm$ 0.36 & 3.55 $\pm$ 0.32 \\
Baseline $+B+E+S+D$ & 5.82 $\pm$ 0.27 & 3.79 $\pm$ 0.54 \\
Baseline $+B+S+D$ & 6.11 $\pm$ 0.34 & 3.55 $\pm$ 0.75 \\
Baseline $+B+D$ & 6.12 $\pm$ 0.35 & 3.70 $\pm$ 0.92 \\
Baseline $+S+D$ & 6.99 $\pm$ 0.47 & 2.93 $\pm$ 0.75 \\
Baseline $+D$ & 6.86 $\pm$ 0.57 & 3.21 $\pm$ 0.68 \\
Baseline $+S$ & 7.64 $\pm$ 0.58 & 2.52 $\pm$ 0.79 \\
Baseline $+E+D$ & 6.01 $\pm$ 0.43 & 4.16 $\pm$ 1.20 \\
Baseline $+B+E+D$ & 5.77 $\pm$ 0.29 & 4.46 $\pm$ 0.93 \\
Baseline $+B$ & 6.32 $\pm$ 0.47 & 4.09 $\pm$ 1.16 \\
Baseline & 7.78 $\pm$ 0.54 & 2.85 $\pm$ 0.82 \\
Baseline $+B+E$ & 5.81 $\pm$ 0.17 & 4.84 $\pm$ 1.10 \\
\end{longtable}
\endgroup
\FloatBarrier
\subsection{Complementary Dice Measures}

\label{sec:metric-sensitivity-details}

The primary Dice scores emphasize the labelled target structures.
Evaluating background separately complements them by checking whether
foreground gains accompany improved agreement in the surrounding region.
The recovered per-volume confusion matrices permit both measurements
on the same saved predictions, without further training or inference.
Write $D_{i,\mathrm{fg}}$ for the primary volume score and
$D_{i,\mathrm{bg}}=D_{i,0}$ for background Dice. Equal weighting of
all classes gives
\begin{equation}
D_{i,\mathrm{all}}=
\begin{cases}
\displaystyle\frac{2D_{i,\mathrm{fg}}+D_{i,\mathrm{bg}}}{3},
 & \text{MSD},\\[.6em]
\displaystyle\frac{D_{i,\mathrm{fg}}+D_{i,\mathrm{bg}}}{2},
 & \text{ADNI and GeoFWI3D}.
\end{cases}
\label{eq:reported-all-class-dice}
\end{equation}
MSD retains separate anterior and posterior classes. The class scores
are calculated within each volume, followed by equal-volume, seed and
fold reductions as in Section~\ref{sec:evaluation-methods}; confusion
counts are not pooled across volumes to form a substitute Dice score.

The complete comparison covers all 504 conventional runs on each
dataset, 108 MSD MedSAM3 run records, all 126 ADNI MedSAM3 runs
and all 126 SAM3 runs.
The MSD $B$ and $B+E$ variants and their matched controls retain complete
foreground results for all 27 runs, but their per-volume records were
not located. Their background and all-class Dice therefore remain
unreported; other baseline records are not substituted for these controls.

Table~\ref{tab:dice-three-measures-summary} illustrates the three measures
for adaptation. The full comparisons are in Appendix Tables~\ref{tab:background-conventional}
and~\ref{tab:background-adaptation}. The conventional and MedSAM3 comparisons retain the
configurations already selected by foreground Dice, allowing the
effect of changing the evaluation measure to be examined directly.


\begin{table}[htbp]
\centering\small\setstretch{1.1}
\begin{tabular}{@{}llrrr@{}}
\hline
Study & Training & Foreground & Background & All classes \\
\hline
SAM3 / GeoFWI3D & Baseline & 69.31 & 97.62 & 83.46 \\
 & $+B+E+D$ & 72.76 & 97.84 & 85.30 \\
MedSAM3 / MSD & Matched baseline & 76.34 & 99.01 & 83.90 \\
 & $+B+E+S+D$ & 81.73 & 99.17 & 87.54 \\
MedSAM3 / ADNI & Baseline & 77.32 & 99.10 & 88.21 \\
 & $+E+S+D$ & 80.85 & 99.19 & 90.02 \\
\hline
\end{tabular}
\caption{Three views of the same predictions: mean Dice (\%) for each
adaptation baseline and the objective selected by foreground Dice.
All-class Dice gives equal weight to each class. Full conventional and
adaptation comparisons are retained in Tables~\ref{tab:background-conventional}
and~\ref{tab:background-adaptation}.}
\label{tab:dice-three-measures-summary}
\end{table}

Every one of the fifteen foreground-selected configurations improves both
mean foreground and mean background Dice. Background means in their
baseline and selected rows range from 95.24\% to 99.76\%, leaving
less room for improvement than the foreground scores. For example,
ADNI MedSAM3 changes from 77.32\% to 80.85\% foreground Dice and
from 99.10\% to 99.19\% background Dice. Its all-class score rises
from 88.21\% to 90.02\%. This supports an improvement in agreement
across both regions, while the foreground measure retains sensitivity
to the target structure.

The higher absolute all-class scores result from the added contribution
of background. Their gains are smaller because background changes
relatively little: its inclusion reduces the foreground weight to
two-thirds on MSD and one-half on the binary tasks. The selected
MSD SwinUNETR gain changes from 11.97 to 8.27 points; the ADNI
MedSAM3 gain changes from 3.53 to 1.81 points. Both quantities describe
the same predictions under different class weights.

The conclusions also hold beyond the selected configurations. Across
all 193 complete constrained-objective comparisons with recoverable
background scores, the same 179 mean effects are positive and the
same fourteen are negative under foreground and all-class Dice.
The highest-mean objective is unchanged in all fifteen model/dataset
comparisons. The selected configurations retain their positive-pair
counts: nine of nine in thirteen comparisons and eight of nine for
geological UNETR and SAM3. Bands alone still improve all 72 conventional medical
model/fold/seed cells under both measures.

This agreement does not imply identical ordering of every outcome.
Six of 1,737 individual paired differences change sign, all on
GeoFWI3D; their foreground magnitudes are below 0.14 points and
their all-class magnitudes below 0.03 points. Some secondary objective
rankings also change. These small differences leave the reported
mean benefits and negative findings intact, without establishing
statistical superiority between closely ranked objectives.

For SAM3, the selected $B+E+D$ variant raises background Dice
from 97.62\% to 97.84\% and all-class Dice from 83.46\% to
85.30\%. Its 3.45-point foreground gain becomes a 1.84-point
all-class gain. All thirteen SAM3 mean effects retain their sign,
and no individual SAM3 paired effect changes sign when background
is included. In particular, the negative $B+E+S$ finding persists
under both measures. These checks use the complete campaign.

\FloatBarrier

\begin{table}[p]
\centering\small\setstretch{1.05}
\setlength{\tabcolsep}{5pt}
\renewcommand{\arraystretch}{1.12}
\begin{tabular}{@{}llrrrrr@{}}
\hline
Model & Variant & $D_{\mathrm{fg}}$ & $D_{\mathrm{bg}}$ & $D_{\mathrm{all}}$ & $\Delta_{\mathrm{fg}}$ & $\Delta_{\mathrm{all}}$ \\
\hline
\multicolumn{7}{@{}l}{\textit{GeoFWI3D}}\\[.25em]
nnU-Net & Baseline & 72.08 & 97.71 & 84.89 & --- & --- \\
 & Baseline $+B+E+S+D$ & 78.22 & 98.29 & 88.25 & +6.14 & +3.36 \\
SwinUNETR & Baseline & 55.54 & 95.24 & 75.39 & --- & --- \\
 & Baseline $+B+E+S+D$ & 57.61 & 95.66 & 76.63 & +2.07 & +1.24 \\
UNETR & Baseline & 54.91 & 96.33 & 75.62 & --- & --- \\
 & Baseline $+B+E$ & 56.07 & 96.39 & 76.23 & +1.16 & +0.61 \\
3D U-Net & Baseline & 77.60 & 98.25 & 87.92 & --- & --- \\
 & Baseline $+B+E+S+D$ & 80.73 & 98.58 & 89.66 & +3.14 & +1.73 \\
\hline
\multicolumn{7}{@{}l}{\textit{MSD}}\\[.25em]
nnU-Net & Baseline & 62.52 & 99.70 & 74.92 & --- & --- \\
 & Baseline $+B+S+D$ & 72.67 & 99.76 & 81.70 & +10.15 & +6.79 \\
SwinUNETR & Baseline & 54.04 & 98.52 & 68.87 & --- & --- \\
 & Baseline $+B+E+S+D$ & 66.00 & 99.40 & 77.13 & +11.97 & +8.27 \\
UNETR & Baseline & 54.71 & 98.96 & 69.46 & --- & --- \\
 & Baseline $+B+E+S$ & 60.07 & 99.18 & 73.11 & +5.35 & +3.64 \\
3D U-Net & Baseline & 73.83 & 99.61 & 82.42 & --- & --- \\
 & Baseline $+B+E+D$ & 80.00 & 99.75 & 86.58 & +6.17 & +4.16 \\
\hline
\multicolumn{7}{@{}l}{\textit{ADNI}}\\[.25em]
nnU-Net & Baseline & 75.66 & 99.74 & 87.70 & --- & --- \\
 & Baseline $+B+S+D$ & 76.69 & 99.76 & 88.23 & +1.03 & +0.52 \\
SwinUNETR & Baseline & 68.72 & 99.57 & 84.14 & --- & --- \\
 & Baseline $+B+E+D$ & 73.03 & 99.66 & 86.34 & +4.31 & +2.20 \\
UNETR & Baseline & 56.27 & 99.24 & 77.76 & --- & --- \\
 & Baseline $+B+E+S+D$ & 60.46 & 99.41 & 79.93 & +4.20 & +2.18 \\
3D U-Net & Baseline & 69.64 & 99.58 & 84.61 & --- & --- \\
 & Baseline $+B+E+D$ & 75.18 & 99.70 & 87.44 & +5.54 & +2.83 \\
\hline
\end{tabular}
\caption{Foreground, background and all-class Dice for the conventional
models. Scores are percentages; paired gains are percentage points.
Each baseline is followed by the configuration selected by its highest
foreground mean in the preceding results. Each row averages three seeds
within each of three folds and then weights folds equally. The selected
configuration is also highest under all-class Dice in every model/dataset
comparison. These tables show means; fold variation and the original
selection limitations remain as reported above. Gains use unrounded values.}
\label{tab:background-conventional}
\end{table}
\FloatBarrier
\begin{table}[p]
\centering\small\setstretch{1.05}
\setlength{\tabcolsep}{5pt}
\renewcommand{\arraystretch}{1.12}
\begin{tabular}{@{}llrrrrr@{}}
\hline
Model & Variant & $D_{\mathrm{fg}}$ & $D_{\mathrm{bg}}$ & $D_{\mathrm{all}}$ & $\Delta_{\mathrm{fg}}$ & $\Delta_{\mathrm{all}}$ \\
\hline
\multicolumn{7}{@{}l}{\textit{MSD: MedSAM3}}\\[.25em]
MedSAM3 & Baseline & 76.34 & 99.01 & 83.90 & --- & --- \\
 & Baseline $+B+E+S+D$ & 81.73 & 99.17 & 87.54 & +5.38 & +3.64 \\
\hline
\multicolumn{7}{@{}l}{\textit{ADNI: MedSAM3}}\\[.25em]
MedSAM3 & Baseline & 77.32 & 99.10 & 88.21 & --- & --- \\
 & Baseline $+E+S+D$ & 80.85 & 99.19 & 90.02 & +3.53 & +1.81 \\
\hline
\multicolumn{7}{@{}l}{\textit{GeoFWI3D: SAM3, all fourteen objectives}}\\[.25em]
SAM3 & Baseline & 69.31 & 97.62 & 83.46 & --- & --- \\
 & Baseline $+B$ & 70.79 & 97.72 & 84.25 & +1.48 & +0.79 \\
 & Baseline $+S$ & 69.10 & 97.62 & 83.36 & -0.22 & -0.11 \\
 & Baseline $+B+S$ & 69.60 & 97.60 & 83.60 & +0.29 & +0.14 \\
 & Baseline $+B+E$ & 69.44 & 97.63 & 83.53 & +0.12 & +0.07 \\
 & Baseline $+B+E+S$ & 67.46 & 97.49 & 82.48 & -1.85 & -0.99 \\
 & Baseline $+D$ & 72.29 & 97.79 & 85.04 & +2.98 & +1.58 \\
 & Baseline $+B+D$ & 72.45 & 97.82 & 85.13 & +3.13 & +1.67 \\
 & Baseline $+E+D$ & 69.18 & 97.61 & 83.40 & -0.13 & -0.07 \\
 & Baseline $+S+D$ & 69.93 & 97.64 & 83.78 & +0.61 & +0.32 \\
 & Baseline $+B+E+D$ & 72.76 & 97.84 & 85.30 & +3.45 & +1.84 \\
 & Baseline $+B+S+D$ & 72.17 & 97.79 & 84.98 & +2.85 & +1.52 \\
 & Baseline $+E+S+D$ & 69.70 & 97.65 & 83.68 & +0.39 & +0.21 \\
 & Baseline $+B+E+S+D$ & 70.95 & 97.73 & 84.34 & +1.64 & +0.88 \\
\hline
\end{tabular}
\caption{Complementary Dice measures for pretrained adaptation. MedSAM3
rows compare the matched baseline with the previously selected
configuration; the MSD row uses the control matched to $B+E+S+D$.
The SAM3 block includes all fourteen completed objective conditions,
including negative effects; $B+E+D$ has the highest foreground and
all-class mean. All rows use three folds and three seeds, the same
per-volume averaging as Table~\ref{tab:background-conventional},
percentage scores and percentage-point gains.}
\label{tab:background-adaptation}
\end{table}
\FloatBarrier
\subsection{Qualitative Prediction Examples}
\label{sec:qualitative-predictions}

The prediction panels supplement the aggregate comparisons with examples
of where models agree with or depart from the reference. They use fresh
matched reproductions, because the original final adapters and prediction
volumes were not retained. For each pretrained-model study, the constrained
objective is the one with the highest original mean validation Dice.
The reproduced run pair was chosen before execution as the pair with the
median original gain among the nine recorded comparisons. Thus, each
dataset contributes one baseline and one constrained reproduction, rather
than a new three-fold, three-seed study.

The repeats preserve the recorded support and validation cases,
initial weights, sampling sequence and update schedule. SAM3 uses five
GeoFWI3D support volumes and 25 updates; MedSAM3 uses five MSD volumes or
two ADNI crops and 200 updates. Medical runs reuse the saved calibration
coefficients. The geological runs repeat the original five-update
calibration procedure. Inference uses images and the fixed text queries;
reference masks enter evaluation and figure selection, not the model's
inference prompts. Case identities, reference hashes and final-adapter
identities were checked, and Dice and error counts were recalculated from
all 416 saved support and validation predictions.

\begin{table}[htbp]
\centering\small\setstretch{1.05}
\renewcommand{\arraystretch}{1.2}
\begin{tabular}{@{}lrrrrr@{}}
\hline
& \multicolumn{2}{c}{Original run pair} &
\multicolumn{3}{c}{Qualitative reproduction} \\
Study (held-out cases) & Baseline & Constrained & Baseline & Constrained & $\Delta$ \\
\hline
GeoFWI3D / SAM3 (96) & 75.39 & 78.21 & 76.72 & 79.47 & +2.75 \\
MSD / MedSAM3 (46) & 77.55 & 82.47 & 78.63 & 81.87 & +3.24 \\
ADNI / MedSAM3 (54) & 78.02 & 81.57 & 77.12 & 80.64 & +3.52 \\
\hline
\end{tabular}
\caption{Mean foreground Dice for the selected original run pairs and
their qualitative reproductions. Scores are percentages and $\Delta$ is
the reproduced paired gain in percentage points, calculated before
rounding. These are single-run validation means, not the original
three-fold aggregates. Constrained objectives follow
Figure~\ref{fig:qualitative-typical}.}
\label{tab:qualitative-reproduction-scores}
\end{table}

The repetitions are not numerically identical to the original runs
(Table~\ref{tab:qualitative-reproduction-scores}). The largest change in
paired gain occurs on MSD: 4.93 points in the selected original pair
becomes 3.24 points. Training records already diverge at the first update,
before inference, despite matching initial weights and sampling. On
GeoFWI3D and ADNI the first forward losses agree but gradient norms differ;
on MSD the baseline forward losses also differ. The recorded MSD hardware
changes from a full A100-SXM4 GPU to an A100 PCIe GPU partition. These
checks locate the divergence in optimization but do not establish its
exact numerical cause. Repeating geological calibration also changes
the edge and depth coefficients. The panels therefore describe the
reproduced models and do not stand in for the original predictions.

For each reproduced validation set, cases are ordered by constrained
minus baseline whole-volume Dice, with case identifier breaking ties.
The upper-median case is shown in Figure~\ref{fig:qualitative-typical};
the minimum and maximum changes are shown below. The same predetermined
slice rule applies to all examples: choose the first slice along source
array axis two, indexed from zero, with the largest reference foreground
area. Images, views, intensity limits and label colours are identical
within each pair. These are stored-grid views, with no anatomical-plane
claim. Scores describe entire volumes, so the displayed slice need not
contain the errors responsible for the largest part of a score change.

The constrained reproduction improves Dice in 88 of 96 GeoFWI3D volumes,
41 of 46 MSD crops and 51 of 54 ADNI crops. The exceptions make the limits
of the average improvement concrete. In the adverse geological case,
whole-volume false positives rise from 4,981 to 7,947, while false
negatives fall only from 911 to 773. The selected slice lies on the cube
face and shows only part of that error balance. In the adverse MSD case,
false positives rise from 188 to 645 despite fewer missed voxels and
anterior/posterior swaps. In the adverse ADNI case, both missed and extra
foreground increase. These examples support the need to examine the
kind and location of errors alongside Dice; they provide no assessment
of clinical or geological usefulness.

\begin{figure}[p]
\centering
\includegraphics[width=\linewidth]{figures/qualitative/geofwi3d-extreme-predictions.png}
\caption{SAM3 on GeoFWI3D: the smallest and largest whole-volume Dice
changes in the reproduced validation set. Blue denotes salt. The adverse
case loses 8.47 points and the largest-gain case gains 7.41 points.
The upper view is slice 95, the last slice of the stored cube. Both views
follow the same reference-area selection rule as the typical examples;
neither was chosen by inspecting prediction errors.}
\label{fig:qualitative-geo-extremes}
\end{figure}

\begin{figure}[p]
\centering
\includegraphics[width=\linewidth]{figures/qualitative/msd-extreme-predictions.png}
\caption{MedSAM3 on MSD: the smallest and largest whole-volume mean
anterior/posterior Dice changes, $-3.02$ and $+9.05$ points. Blue denotes
anterior and orange posterior hippocampus. The adverse example illustrates
that extending a prediction can introduce disagreement even when more
reference foreground is recovered across the volume.}
\label{fig:qualitative-msd-extremes}
\end{figure}

\begin{figure}[p]
\centering
\includegraphics[width=\linewidth]{figures/qualitative/adni-extreme-predictions.png}
\caption{MedSAM3 on ADNI: the smallest and largest whole-volume Dice
changes, $-2.19$ and $+11.63$ points. Blue denotes hippocampus. In the
largest-gain example, the constrained model retains more of the elongated
reference region, although substantial disagreement remains. Each row
uses one common image and slice for both models.}
\label{fig:qualitative-adni-extremes}
\end{figure}

\FloatBarrier
\subsection{Annotation Compatibility and the Anterior/Posterior Transition}
\label{sec:annotation-results}

The audit in Section~\ref{sec:annotation-compatibility} examines whether
the published anterior/posterior annotation convention can support an
additional anatomical constraint. It tests compatibility with the
released labels and the feasibility of locating the transition
automatically.

The MSD dataset description attributes manual tracing to two earlier
hippocampal annotation protocols \cite{simpson2019dataset}. Pruessner
\emph{et al.}\ \cite{pruessner2000volumetry} describe manual volumetry with
simultaneous inspection of sagittal, coronal and horizontal views; the last
correspond to planes dividing upper and lower portions. Woolard and Heckers
\cite{woolard2012asymmetry} build on that protocol: sagittal tracing is
refined in coronal views, and the hippocampus is then divided into anterior
uncus and posterior body and tail. Their anatomical cues include ``more than
one cut through the hippocampus'' in a coronal section and a protrusion
towards the brain's midline and the crus cerebri, a nearby midbrain
structure; the division is checked in the sagittal view.

The MSD paper specifies the final head section as the ``coronal slice
containing the uncal apex'' \cite{simpson2019dataset}. These descriptions
motivate two candidate constraints: a single coronal A/P separation and
a transition located from the appearance or disappearance of the folded
uncal profile. The audit asks whether these descriptions yield
reproducible rules for the released crops. Exact separation on the stored
voxel grid and disconnected mask components are operational proxies;
the papers describe an anatomical tracing procedure using multiple
views, not these numerical tests.

The planar analysis finds that 58 of 260 crops cannot be reproduced
exactly by a single coronal A/P division. The best planes disagree
with 1,388 of 856,754 foreground voxels overall, approximately
0.162\%. A strict universal rule therefore has counterexamples,
although the aggregate discrepancy is small. This supports treating
the division as broadly plane-like while retaining its actual discrete
annotation for supervision.

% PROFILE AUDIT: experiments/uncal_validation_profiles_20260921/audit.json;
% metric_counts: split_min2_at_cut=2, split_min2_near_cut=2, n=52.
The profile analysis examines a numerical proxy for the multiple-cut
cue described by Woolard and Heckers \cite{woolard2012asymmetry}.
Only two of 52 validation crops contain two disconnected foreground
components of at least two voxels each in the final anterior section;
allowing a window of three sections on either side of the fitted
interface still identifies only two crops. This criterion therefore
does not recover the annotated transition consistently. A folded
three-dimensional structure can intersect one plane several times,
and distinct tissue profiles can remain connected in the discrete mask.
Nor does the absence of a visible separation prove the absence of the
anatomical landmark. The explored image cues did not establish a
universally reliable automatic transition location.

\begin{table}[htbp]
\centering\small\setstretch{1.08}
\begin{tabular}{@{}P{3.6cm}P{4.2cm}P{6.8cm}@{}}
\hline
Candidate rule & Observed check & Consequence for formulation \\
\hline
Single coronal A/P plane, motivated by the MSD convention & 58/260 crops not exactly planar; 0.162\% aggregate foreground disagreement & A strict universal plane does not reproduce all stored labels; the voxel-level discrepancy is small. \\[.6em]
Disconnected profiles as a proxy for the Woolard--Heckers cue & 2/52 crops in the last anterior section; 2/52 within three sections of the fitted interface & The component test does not consistently identify the annotated transition or independently locate the uncal apex. \\
\hline
\end{tabular}
\caption{Compatibility of two proposed A/P transition rules with the
released MSD labels. The tests are motivated by the MSD descriptor and
its cited annotation protocol \cite{simpson2019dataset,woolard2012asymmetry}.}
\label{tab:annotation-audit-results}
\end{table}

These findings guide the constraint design: the annotated A/P interface
provides the classwise supervision, while a strict plane or an automatic
uncal-profile rule is not imposed universally. This avoids penalizing
predictions that agree with the stored labels but violate a simplified
geometric proxy. The audit establishes the behaviour of those proxies;
without an expert landmark reference, it does not determine whether
the original annotators misapplied the anatomical protocol.


\FloatBarrier

\subsection{Geological Development Comparison}\label{sec:geo-development}
An earlier UNETR development comparison motivated the observation-frame
exclusion described in Section~\ref{sec:geo-methods}. At learning rate
$5\times10^{-4}$, on fold 0 and seed 0, excluding the frame and
recalibrating the band coefficient increased final validation Dice
from 50.61\% to 52.03\% relative to the previous band implementation
(+1.42 percentage points); both runs used 407 updates. This supported
retaining the modification as part of the generalization strategy.
The comparison concerns one development pair and the joint change in
support and calibration; it does not establish a universal benefit
from frame exclusion. It is separate from the final three-fold,
three-seed matrix, whose UNETR recipe uses learning rate
$6\times10^{-4}$.
% Development evidence: reports/geofwi3d_unetr_lr5_frame_bands_20261019/
% observed_final_status.json and protocol_difference.json.


\subsection{Geological Comparison}

\begin{table}[htbp]
\centering\small\setstretch{1.05}
\setlength{\tabcolsep}{4pt}
\renewcommand{\arraystretch}{1.12}
\begin{tabular}{@{}l*{4}{rr}@{}}
\hline
 & \multicolumn{2}{c}{nnU-Net} & \multicolumn{2}{c}{SwinUNETR} & \multicolumn{2}{c}{UNETR} & \multicolumn{2}{c}{3D U-Net} \\
Training variant & $\Delta$ & $n_+$ & $\Delta$ & $n_+$ & $\Delta$ & $n_+$ & $\Delta$ & $n_+$ \\
\hline
Baseline $+B$ & +1.35 & 8/9 & +1.59 & 8/9 & +0.75 & 7/9 & +0.72 & 7/9 \\
Baseline $+S$ & -0.07 & 4/9 & +0.37 & 7/9 & -0.33 & 2/9 & +0.24 & 7/9 \\
Baseline $+B+S$ & +1.06 & 6/9 & +1.44 & 8/9 & +0.06 & 3/9 & +0.31 & 7/9 \\
Baseline $+B+E$ & +0.76 & 7/9 & +1.45 & 9/9 & +1.16 & 8/9 & +0.55 & 7/9 \\
Baseline $+B+E+S$ & -0.29 & 5/9 & +1.39 & 8/9 & -0.20 & 3/9 & +0.96 & 7/9 \\
Baseline $+D$ & +1.84 & 9/9 & +1.01 & 9/9 & +0.68 & 6/9 & +1.81 & 9/9 \\
Baseline $+B+D$ & +3.15 & 9/9 & +1.74 & 9/9 & +0.78 & 6/9 & +2.80 & 9/9 \\
Baseline $+E+D$ & +3.13 & 9/9 & +1.21 & 9/9 & +0.77 & 6/9 & +1.99 & 8/9 \\
Baseline $+S+D$ & +4.18 & 9/9 & +0.96 & 8/9 & +0.21 & 6/9 & +2.17 & 9/9 \\
Baseline $+B+E+D$ & -1.36 & 5/9 & +1.79 & 9/9 & +0.94 & 7/9 & +2.90 & 9/9 \\
Baseline $+B+S+D$ & +5.25 & 9/9 & +2.04 & 9/9 & +0.41 & 5/9 & +3.13 & 9/9 \\
Baseline $+E+S+D$ & +4.77 & 8/9 & +1.22 & 9/9 & -0.24 & 4/9 & +2.69 & 9/9 \\
Baseline $+B+E+S+D$ & +6.14 & 9/9 & +2.07 & 9/9 & +0.65 & 6/9 & +3.14 & 9/9 \\
\hline
\end{tabular}
\caption{GeoFWI3D: mean paired validation-Dice gains $\Delta$
(percentage points) and positive fold/seed pairs $n_+$ out of nine,
for every constrained objective. Gains use unrounded matched scores.}
\label{tab:geo-ledger}
\end{table}

\FloatBarrier
\subsection{Conventional Medical Comparisons}

\begin{table}[htbp]
\centering\small\setstretch{1.05}
\setlength{\tabcolsep}{4pt}
\renewcommand{\arraystretch}{1.12}
\begin{tabular}{@{}l*{4}{rr}@{}}
\hline
 & \multicolumn{2}{c}{nnU-Net} & \multicolumn{2}{c}{SwinUNETR} & \multicolumn{2}{c}{UNETR} & \multicolumn{2}{c}{3D U-Net} \\
Training variant & $\Delta$ & $n_+$ & $\Delta$ & $n_+$ & $\Delta$ & $n_+$ & $\Delta$ & $n_+$ \\
\hline
Baseline $+B$ & +7.55 & 9/9 & +3.85 & 9/9 & +3.89 & 9/9 & +4.83 & 9/9 \\
Baseline $+S$ & -2.04 & 0/9 & +0.26 & 6/9 & +0.78 & 6/9 & -0.17 & 5/9 \\
Baseline $+B+S$ & +7.59 & 9/9 & +4.97 & 9/9 & +5.14 & 9/9 & +4.85 & 9/9 \\
Baseline $+B+E$ & +6.45 & 9/9 & +4.00 & 9/9 & +4.13 & 9/9 & +5.71 & 9/9 \\
Baseline $+B+E+S$ & +6.47 & 9/9 & +5.16 & 9/9 & +5.35 & 9/9 & +5.53 & 9/9 \\
Baseline $+D$ & +1.61 & 8/9 & +8.24 & 9/9 & +1.76 & 8/9 & +2.30 & 9/9 \\
Baseline $+B+D$ & +9.38 & 9/9 & +11.04 & 7/9 & +3.70 & 9/9 & +5.02 & 9/9 \\
Baseline $+E+D$ & +0.46 & 6/9 & +8.97 & 9/9 & +2.21 & 9/9 & +3.19 & 8/9 \\
Baseline $+S+D$ & +3.00 & 9/9 & +8.44 & 9/9 & +2.37 & 9/9 & +2.36 & 9/9 \\
Baseline $+B+E+D$ & +9.09 & 9/9 & +11.55 & 7/9 & +3.80 & 9/9 & +6.17 & 9/9 \\
Baseline $+B+S+D$ & +10.15 & 9/9 & +11.67 & 9/9 & +4.70 & 9/9 & +4.95 & 9/9 \\
Baseline $+E+S+D$ & +2.39 & 8/9 & +9.12 & 9/9 & +2.71 & 9/9 & +3.33 & 8/9 \\
Baseline $+B+E+S+D$ & +9.98 & 9/9 & +11.97 & 9/9 & +4.89 & 9/9 & +6.07 & 9/9 \\
\hline
\end{tabular}
\caption{MSD: mean paired validation-Dice gains $\Delta$
(percentage points) and positive fold/seed pairs $n_+$ out of nine,
for every constrained objective. Gains use unrounded matched scores.}
\label{tab:msd-noaug-ledger}
\end{table}

\begin{table}[htbp]
\centering\small\setstretch{1.05}
\setlength{\tabcolsep}{4pt}
\renewcommand{\arraystretch}{1.12}
\begin{tabular}{@{}l*{4}{rr}@{}}
\hline
 & \multicolumn{2}{c}{nnU-Net} & \multicolumn{2}{c}{SwinUNETR} & \multicolumn{2}{c}{UNETR} & \multicolumn{2}{c}{3D U-Net} \\
Training variant & $\Delta$ & $n_+$ & $\Delta$ & $n_+$ & $\Delta$ & $n_+$ & $\Delta$ & $n_+$ \\
\hline
Baseline $+B$ & +0.73 & 9/9 & +3.34 & 9/9 & +2.99 & 9/9 & +3.81 & 9/9 \\
Baseline $+S$ & +0.01 & 6/9 & -0.07 & 3/9 & -0.34 & 1/9 & -0.24 & 2/9 \\
Baseline $+B+S$ & +0.73 & 9/9 & +3.27 & 9/9 & +3.24 & 9/9 & +4.67 & 9/9 \\
Baseline $+B+E$ & +0.59 & 8/9 & +3.68 & 9/9 & +3.49 & 9/9 & +4.57 & 9/9 \\
Baseline $+B+E+S$ & +0.60 & 8/9 & +3.62 & 9/9 & +3.91 & 9/9 & +4.76 & 9/9 \\
Baseline $+D$ & +0.31 & 9/9 & +1.35 & 9/9 & +1.51 & 9/9 & +1.97 & 8/9 \\
Baseline $+B+D$ & +0.96 & 9/9 & +3.91 & 9/9 & +3.70 & 9/9 & +4.51 & 9/9 \\
Baseline $+E+D$ & +0.48 & 9/9 & +1.78 & 9/9 & +1.85 & 9/9 & +2.30 & 9/9 \\
Baseline $+S+D$ & +0.43 & 9/9 & +1.31 & 9/9 & +1.52 & 9/9 & +1.98 & 9/9 \\
Baseline $+B+E+D$ & +0.87 & 9/9 & +4.31 & 9/9 & +4.10 & 9/9 & +5.54 & 9/9 \\
Baseline $+B+S+D$ & +1.03 & 9/9 & +3.80 & 9/9 & +3.52 & 9/9 & +4.72 & 9/9 \\
Baseline $+E+S+D$ & +0.53 & 8/9 & +1.76 & 9/9 & +2.07 & 9/9 & +2.75 & 9/9 \\
Baseline $+B+E+S+D$ & +0.88 & 9/9 & +4.22 & 9/9 & +4.20 & 9/9 & +5.39 & 9/9 \\
\hline
\end{tabular}
\caption{ADNI: mean paired validation-Dice gains $\Delta$
(percentage points) and positive fold/seed pairs $n_+$ out of nine,
for every constrained objective. Gains use unrounded matched scores.}
\label{tab:adni-noaug-ledger}
\end{table}

\FloatBarrier

\subsection{SAM3 Fold Effects}

\begin{table}[htbp]
\centering\small\setstretch{1.05}
\setlength{\tabcolsep}{5pt}
\renewcommand{\arraystretch}{1.12}
\begin{tabular}{@{}lrrrrrr@{}}
\hline
 & \multicolumn{3}{c}{Validation Dice (\%)} & \multicolumn{3}{c}{$\Delta$ Dice (pp)} \\
Variant & Fold 0 & Fold 1 & Fold 3 & Fold 0 & Fold 1 & Fold 3 \\
\hline
Baseline & 64.77 & 72.30 & 70.87 & 0.00 & 0.00 & 0.00 \\
Baseline $+B$ & 70.65 & 73.23 & 68.48 & +5.88 & +0.94 & -2.39 \\
Baseline $+S$ & 65.04 & 74.51 & 67.74 & +0.27 & +2.21 & -3.13 \\
Baseline $+B+S$ & 69.29 & 71.28 & 68.25 & +4.52 & -1.02 & -2.62 \\
Baseline $+B+E$ & 65.45 & 71.46 & 71.40 & +0.68 & -0.84 & +0.53 \\
Baseline $+B+E+S$ & 66.01 & 68.11 & 68.27 & +1.24 & -4.19 & -2.61 \\
Baseline $+D$ & 70.12 & 74.12 & 72.64 & +5.35 & +1.83 & +1.77 \\
Baseline $+B+D$ & 71.66 & 73.02 & 72.66 & +6.89 & +0.72 & +1.79 \\
Baseline $+E+D$ & 65.85 & 70.11 & 71.58 & +1.08 & -2.19 & +0.71 \\
Baseline $+S+D$ & 69.85 & 71.25 & 68.67 & +5.08 & -1.05 & -2.20 \\
Baseline $+B+E+D$ & 70.97 & 74.87 & 72.45 & +6.19 & +2.57 & +1.58 \\
Baseline $+B+S+D$ & 71.14 & 72.68 & 72.67 & +6.37 & +0.39 & +1.80 \\
Baseline $+E+S+D$ & 67.81 & 67.86 & 73.42 & +3.04 & -4.44 & +2.55 \\
Baseline $+B+E+S+D$ & 69.30 & 71.94 & 71.62 & +4.52 & -0.35 & +0.75 \\
\hline
\end{tabular}
\caption{SAM3 fold means, each averaging seeds 0, 1 and 2. Paired changes
use the corresponding baseline within each fold. All 126 final runs
are included, and every objective is shown regardless of the direction
of its change.}
\label{tab:sam3-geo-folds}
\end{table}

\FloatBarrier

\subsection{MedSAM3 Fold Effects}

\begin{table}[htbp]
\centering\small\setstretch{1.05}
\setlength{\tabcolsep}{4pt}
\renewcommand{\arraystretch}{1.12}
\begin{tabular}{@{}lrrrrrr@{}}
\hline
 & \multicolumn{3}{c}{Validation Dice (\%)} & \multicolumn{3}{c}{$\Delta$ Dice (pp)} \\
Training variant & Fold 0 & Fold 1 & Fold 3 & Fold 0 & Fold 1 & Fold 3 \\
\hline
\multicolumn{7}{@{}l}{\textbf{MSD}}\\[.3em]
Baseline for $B$, $B+E$ & 73.46 & 73.89 & 77.43 & +0.00 & +0.00 & +0.00 \\
Baseline $+B$ & 78.36 & 78.89 & 79.03 & +4.89 & +5.00 & +1.60 \\
Baseline $+B+E$ & 79.60 & 79.26 & 81.50 & +6.14 & +5.38 & +4.07 \\
Baseline for other variants & 77.35 & 74.45 & 77.23 & 0.00 & 0.00 & 0.00 \\
Baseline $+B+E+S+D$ & 82.08 & 80.16 & 82.94 & +4.74 & +5.71 & +5.71 \\
\hline\multicolumn{7}{@{}l}{\textbf{ADNI}}\\[.3em]
Baseline & 75.21 & 77.70 & 79.06 & 0.00 & 0.00 & 0.00 \\
Baseline $+E+S+D$ & 78.57 & 81.47 & 82.52 & +3.36 & +3.77 & +3.46 \\
\hline
\end{tabular}
\caption{MedSAM3 fold means, each averaging three seeds. The
highest-mean objective for each dataset is selected using the overall
validation mean. MSD also shows $B$ and $B+E$ with their matched control;
the two baseline rows preserve the corresponding execution-batch pairings.
Gains use unrounded paired scores.}
\label{tab:medsam-noaug-folds}
\label{tab:medsam-msd-original-folds}
\end{table}

\clearpage

\section{Reproducibility Material}
\label{sec:evidence-register}
\label{sec:background-dice-evidence}
\label{sec:all-variant-error-evidence}

The project archive links run identifiers to source versions,
configurations, fold and seed inventories, support/evaluation manifests,
calibrated coefficients and final metrics. The directories below, relative
to \path{reports/}, contain the thesis extracts and their verification
records. Their scripts and metadata identify the underlying campaign
archives; full training-set scores and update counts remain available
in the run-level comma-separated values (CSV) and JavaScript Object Notation
(JSON) files.

\begin{center}
\small\setstretch{1.05}
\renewcommand{\arraystretch}{1.2}
\begin{tabular}{@{}P{3.8cm}P{11.2cm}@{}}
\hline
Evidence & Archive entry point \\
\hline
Conventional GeoFWI3D & \path{thesis_geofwi_update_20261022/} \\
MSD/ADNI conventional and MedSAM3 & \path{thesis_noaug_update_20261021/} \\
Complete SAM3 results & \path{thesis_sam3_final_20261023/} \\
Per-volume confusion matrices and Dice & \path{background_dice_archive_recovery_20261023/} \\
Complementary Dice tables & \path{thesis_background_dice_update_20261023/} \\
Full-validation error tables & \path{thesis_all_variant_errors_20261023/} \\
Qualitative reproductions and panels & \path{thesis_qualitative_predictions_20261103/} \\
\hline
\end{tabular}
\end{center}

Aggregates are reconstructed from final run records, averaging seeds
within each fold and then weighting folds equally. Reported dispersion
is the sample standard deviation of fold means. Paired gains retain the
baseline for the same fold, seed and execution batch, including the
distinct matched controls for MSD MedSAM3 $B$ and $B+E$. Aggregate verification
records distinguish inventory checks, numerical recalculation and
available payload-hash checks; these checks use retained records rather
than rerunning training or inference.

The qualitative archive is separate from those aggregate extracts. It
contains the six reproduction records, saved image/reference/prediction
arrays, original-score comparisons, verification receipts and the scripts
and selection metadata for the panels. Section~\ref{sec:qualitative-predictions}
states their provenance and numerical differences; none of these repeats
changes the original aggregate tables or charts.

Saved confusion matrices support recalculation of per-volume Dice and
full-grid errors. The complete SAM3 extraction supplies the final 126
evaluations and replaces the earlier SAM3 snapshots in the complementary
Dice and error extracts. For MSD MedSAM3 $B$, $B+E$ and their control,
the 27 aggregate foreground scores and error records are preserved;
their per-volume matrices were not recovered. Consequently, their
background and all-class scores remain unavailable, as stated in
Section~\ref{sec:background-dice-results}. All six studies nevertheless
have completed results for their declared objective inventories.

Detailed archive inventories, execution histories and verification
receipts remain with the corresponding research records.

\clearpage



